\documentclass[12pt, a4paper]{report}

% -- Encoding & language ------------------------------------------------------
\usepackage[utf8]{inputenc}
\usepackage[T1]{fontenc}
\usepackage[english]{babel}

% -- Math ---------------------------------------------------------------------
\usepackage{amsmath, amssymb, amsthm, mathtools}
\usepackage{bm}

% -- Layout -------------------------------------------------------------------
\usepackage[top=2.5cm, bottom=2.5cm, left=2.8cm, right=2.8cm]{geometry}
\usepackage{parskip}
\setlength{\parskip}{6pt}

% -- Tables -------------------------------------------------------------------
\usepackage{booktabs}
\usepackage{array}
\usepackage{tabularx}

% -- Colors & boxes -----------------------------------------------------------
\usepackage{xcolor}
\usepackage[most]{tcolorbox}
\tcbuselibrary{theorems, skins, breakable}

\definecolor{noteblue}{RGB}{220, 235, 252}
\definecolor{noteblueborder}{RGB}{52, 120, 200}
\definecolor{warnred}{RGB}{252, 228, 220}
\definecolor{warnredborder}{RGB}{200, 72, 52}
\definecolor{defgreen}{RGB}{220, 248, 228}
\definecolor{defgreenborder}{RGB}{40, 160, 80}
\definecolor{propgold}{RGB}{252, 248, 220}
\definecolor{propgoldborder}{RGB}{180, 150, 40}

\newtcolorbox{notebox}[1][Nota]{
  breakable,
  colback=noteblue, colframe=noteblueborder,
  fonttitle=\bfseries\small, title={#1},
  boxrule=0.8pt, arc=4pt,
  left=6pt, right=6pt, top=4pt, bottom=4pt
}

\newtcolorbox{warnbox}[1][Attenzione]{
  breakable,
  colback=warnred, colframe=warnredborder,
  fonttitle=\bfseries\small, title={#1},
  boxrule=0.8pt, arc=4pt,
  left=6pt, right=6pt, top=4pt, bottom=4pt
}

\newtcolorbox{defbox}[1][Definizione]{
  breakable,
  colback=defgreen, colframe=defgreenborder,
  fonttitle=\bfseries\small, title={#1},
  boxrule=0.8pt, arc=4pt,
  left=6pt, right=6pt, top=4pt, bottom=4pt
}

\newtcolorbox{propbox}[1][Proprietà]{
  breakable,
  colback=propgold, colframe=propgoldborder,
  fonttitle=\bfseries\small, title={#1},
  boxrule=0.8pt, arc=4pt,
  left=6pt, right=6pt, top=4pt, bottom=4pt
}

% -- Hyperlinks ---------------------------------------------------------------
\usepackage[colorlinks=true, linkcolor=noteblueborder,
            urlcolor=noteblueborder, citecolor=defgreenborder]{hyperref}

% -- Math operators -----------------------------------------------------------
\DeclareMathOperator{\trace}{trace}
\DeclareMathOperator{\rank}{rank}
\DeclareMathOperator{\sgn}{sgn}
\newcommand{\R}{\mathbb{R}}
\newcommand{\norm}[1]{\left\lVert#1\right\rVert}

% -----------------------------------------------------------------------------
\title{%
  \vspace{-1cm}
  {\LARGE\bfseries Computer Vision}\\[0.6em]
  {\large Note complete del corso}\\[0.3em]
  {\normalsize Capitoli 1\,--\,19}
}
\author{Prof.\ Irene Amerini}
\date{Anno Accademico 2024/2025}
% -----------------------------------------------------------------------------

\begin{document}
\maketitle
\tableofcontents
\newpage

% =============================================================================
\chapter{Introduzione alla Computer Vision}
% =============================================================================

\section{Definizione e obiettivi}

La \textbf{Computer Vision} mira a convertire luce in significato, estraendo informazioni
geometriche, semantiche e temporali da immagini e video.  A differenza
dell'\textit{Image Processing}, che trasforma immagini in altre immagini, la Computer
Vision cerca di inferire proprietà del mondo a partire da dati visivi.

\begin{defbox}[Problema inverso mal posto]
La Computer Vision può essere vista come un \textbf{problema inverso mal posto}: date
una o più immagini 2D, l'obiettivo è recuperare informazioni sulla scena 3D
sottostante.  Molte scene 3D diverse, condizioni di illuminazione e configurazioni
di oggetti possono produrre immagini 2D simili; gli algoritmi necessitano quindi di
assunzioni, vincoli e prior appresi per interpretare correttamente le immagini.
\end{defbox}

\section{Rappresentazione di un'immagine}

Un'immagine può essere interpretata in due modi:

\begin{itemize}
  \item \textbf{Dominio discreto}: griglia di valori di intensità in $[0,\,2^{\text{bit}}]$.
  \item \textbf{Dominio continuo}: una funzione $f:\R^2\to\R$, dove $f(x,y)$ fornisce
        l'intensità nella posizione spaziale $(x,y)$.  Un'immagine digitale è
        semplicemente una versione discreta, campionata e quantizzata di questa funzione
        continua.
\end{itemize}

\section{Computer Vision, Image Processing e Computer Graphics}

\begin{itemize}
  \item \textbf{Image Processing}: trasforma immagini in altre immagini (filtraggio,
        riduzione del rumore, rilevamento di bordi, \ldots).
  \item \textbf{Computer Vision}: inferisce informazioni dal mondo a partire da immagini
        (identità degli oggetti, struttura 3D, moto, posa, semantica della scena).
  \item \textbf{Computer Graphics}: risolve il \emph{problema diretto} --- data una scena
        3D con oggetti, materiali, illuminazione e parametri della telecamera, genera
        un'immagine 2D.  La Computer Vision risolve il \emph{problema inverso}.
\end{itemize}

\section{Percezione visiva}

La percezione visiva \emph{non} è una semplice misurazione delle intensità dei pixel.
La stessa intensità misurata può essere percepita in modo diverso a seconda del
contesto, delle ombre, della geometria e delle conoscenze pregresse.  Ciò spiega
perché i sistemi di visione hanno bisogno di modelli, assunzioni e vincoli per
interpretare correttamente le immagini.

\section{Sfide principali}

Variazione di punto di vista, variazione di scala, deformazione, occlusione,
variazioni di illuminazione, cluttering di sfondo, variazione intra-classe, moto e
ambiguità locale.

\section{Breve contesto storico}

Milestones importanti: stereo e fotogrammetria precoci; shape from shading;
flusso ottico; matrice essenziale; Markov Random Fields; feature locali (SIFT);
reti neurali convoluzionali; ImageNet e AlexNet; metodi di rilevamento di oggetti
(YOLO, Mask R-CNN); modelli di diffusione e modelli multimodali.

\section{Sfide aperte}

Apprendimento non supervisionato e auto-supervisionato, apprendimento interattivo,
robustezza, generalizzazione, bias induttivi, interpretabilità, costi computazionali
e problemi etici/legali.

% =============================================================================
\chapter{Image Processing}
% =============================================================================

\section{Pipeline tipica di acquisizione a colori}

Un'immagine digitale è il risultato di una pipeline di acquisizione ed elaborazione.
La luce viene catturata da un sensore CCD/CMOS dotato di un \textbf{color filter
array} (solitamente un pattern di Bayer).  Le misurazioni grezze del sensore vengono
poi elaborate attraverso: guadagno ISO, elaborazione raw, demosaicing, bilanciamento
del bianco, trasformazione dello spazio colore, riduzione del rumore, mappatura a
sRGB e infine compressione (es.\ JPEG).

\begin{notebox}[Importanza per la Computer Vision]
Le immagini usate dagli algoritmi di CV spesso non sono misurazioni fisiche grezze:
sono già state tone-mappate, corrette per il colore, derumorose e compresse.  Molti
algoritmi assumono che le intensità siano in uno \emph{spazio lineare}.
\end{notebox}

% -----------------------------------------------------------------------------
\section{Operatori puntuali (Point Processing)}
% -----------------------------------------------------------------------------

\begin{defbox}[Operatori puntuali]
Gli operatori puntuali modificano il valore di ciascun pixel \textbf{indipendentemente
dai suoi vicini}.  Il nuovo valore di un pixel in posizione $(x,y)$ dipende solo dal
valore originale in quella posizione.
\end{defbox}

\textbf{Esempi:}
\begin{itemize}
  \item \textbf{Luminosità}: $g(x,y)=f(x,y)+b$.
        Aggiungere una costante $b$ schiarisce o scurisce l'immagine.
  \item \textbf{Contrasto}: $g(x,y)=a\cdot f(x,y)$.
        Moltiplicare per $a>1$ aumenta il contrasto; $a<1$ lo riduce.
  \item \textbf{Inversione}: $g(x,y)=255-f(x,y)$ per un'immagine a 8 bit.
  \item \textbf{Flip orizzontale}: $g(x,y)=f(-x,y)$.
\end{itemize}

% -- 2.1 Gamma correction ------------------------------------------------------
\subsection{Correzione gamma}

La correzione gamma \textbf{aggiusta la luminosità} di un'immagine visualizzata su
schermo per renderla più naturale all'occhio:
\[
  V_{\text{out}} = A \cdot V_{\text{in}}^{\gamma}
\]
dove $V_{\text{in}}$ è la tensione d'ingresso, $V_{\text{out}}$ è l'uscita
regolata, $A$ è una costante e tipicamente $\gamma = 2.2$.

\begin{notebox}[Gamma e spazio sRGB]
La maggior parte delle immagini è memorizzata in spazio sRGB, già tone-mappato e
\emph{non} lineare rispetto alla luminanza della scena.  Molti algoritmi di CV
assumono uno spazio lineare.  In assenza della curva di risposta della camera, si
usa tipicamente $\gamma=2.2$ come approssimazione (non fisicamente esatta, ma spesso
sufficiente).
\end{notebox}

% -- 2.2 Histogram equalization -----------------------------------------------
\subsection{Equalizzazione dell'istogramma}

L'\textbf{Histogram Equalization (HE)} aggiusta il contrasto globale redistribuendo
i valori di intensità per produrre un istogramma più piatto, rendendo i dettagli più
visibili.

\begin{enumerate}
  \item \textbf{Calcolo dell'istogramma}: conta il numero di pixel per ogni livello
        di intensità.
        \[
          p(r_k) = \frac{\#\text{pixel con intensità }r_k}{\#\text{pixel totali}}
        \]

  \item \textbf{Calcolo della CDF}: mappa ogni intensità di input a un valore in
        $[0,1]$.
        \[
          \mathrm{CDF}(r_k) = \sum_{i=0}^{k} p(r_i)
        \]

  \item \textbf{Scalatura}: la CDF viene scalata all'intero range di intensità e
        usata come lookup table.
        \[
          r'_k = \mathrm{round}\!\left(
            \frac{\bigl(\mathrm{CDF}(r_k)-\mathrm{CDF}_{\min}\bigr)(L-1)}
                 {N - \mathrm{CDF}_{\min}}
          \right)
        \]
        dove $L$ è il massimo valore di pixel, $N$ il numero totale di pixel e
        $\mathrm{CDF}_{\min}$ il minimo valore non-zero della CDF.
\end{enumerate}

\textbf{Local Histogram Equalization (AHE):} calcola una trasformazione per ciascun
pixel basandosi sull'istogramma di un piccolo intorno locale, migliorando i dettagli
locali a scapito della complessità computazionale.

\begin{notebox}[CLAHE --- Contrast Limited AHE]
Un limite dell'AHE è che può amplificare fortemente il rumore in regioni quasi
uniformi.  La \textbf{CLAHE} risolve il problema agganciando i picchi dell'istogramma
tramite un \emph{clip limit} predefinito e ridistribuendo i pixel tagliati.
L'interpolazione bilineare tra tile vicine evita artefatti a blocchi.
\end{notebox}

% =============================================================================
\section{Filtraggio lineare}
% =============================================================================

\begin{defbox}[Filtraggio lineare]
Il filtraggio crea una nuova immagine in cui il valore di ciascun pixel è una
\textbf{funzione dei suoi vicini} nell'immagine originale.  Più specificamente, il
\emph{filtraggio lineare} calcola ciascun nuovo pixel come \textbf{combinazione
lineare} (somma pesata) dei suoi vicini; l'insieme dei pesi è detto \textbf{kernel}
o \textbf{filtro}.
\end{defbox}

% -- 2.3.1 Cross-correlation e convoluzione ------------------------------------
\subsection{Cross-correlazione e convoluzione}

Queste sono le due operazioni primarie di filtraggio lineare.

\textbf{Cross-correlazione:} l'output $G[i,j]$ è la somma dei prodotti element-wise
tra il kernel $H$ e il vicinato dell'immagine $F$ centrato in $(i,j)$:
\[
  G[i,j] = \sum_{u=-k}^{k}\sum_{v=-k}^{k} H[u,v]\,F[i+u,\,j+v]
\]

\textbf{Convoluzione:} uguale alla cross-correlazione, tranne che il kernel $H$
viene ribaltato orizzontalmente e verticalmente prima di essere applicato:
\[
  G[i,j] = \sum_{u=-k}^{k}\sum_{v=-k}^{k} H[u,v]\,F[i-u,\,j-v]
\]

\begin{propbox}[Proprietà della convoluzione]
\begin{itemize}
  \item \textbf{Commutativa}: $a*b = b*a$
  \item \textbf{Associativa}: $(a*b_1)*b_2 = a*(b_1*b_2)$.
        Applicare due filtri in sequenza equivale a convolverli prima in un unico
        filtro.
  \item \textbf{Distributiva sull'addizione}.
\end{itemize}
In pratica, se il kernel è simmetrico, convoluzione e cross-correlazione sono
identiche.
\end{propbox}

% -- 2.3.2 Padding e bordi ----------------------------------------------------
\subsection{Padding e gestione dei bordi}

Quando il kernel è centrato su un pixel vicino al bordo dell'immagine, la finestra
del kernel può sconfinare oltre il bordo.  Strategie di padding:

\begin{itemize}
  \item \textbf{Zero}: si aggiungono zeri all'esterno.
  \item \textbf{Costante}: si aggiunge un valore costante $c$.
  \item \textbf{Clamp / replica}: si replica il bordo.
  \item \textbf{Reflection / mirror}: si riflette verso l'interno.
  \item \textbf{Symmetric}: si riflette includendo il bordo.
\end{itemize}

\textbf{Dimensione dell'output.}
Per una convoluzione con input di dimensione $N$, kernel di dimensione $K$, padding
$P$ e stride $S$:
\[
  \dim_{\text{out}} = \left\lfloor\frac{N + 2P - K}{S}\right\rfloor + 1
\]
Casi speciali (stride $S=1$):
\begin{itemize}
  \item \texttt{full}: output più grande dell'input.
  \item \texttt{same}: output con stessa dimensione dell'input
        ($P = \lfloor K/2\rfloor$).
  \item \texttt{valid}: nessun padding, output più piccolo
        ($\dim_{\text{out}} = N - K + 1$).
\end{itemize}

% -- 2.3.3 Esempi di filtri lineari -------------------------------------------
\subsection{Esempi di filtri lineari}

\paragraph{Filtro box (mean filter).}
Sostituisce ogni pixel con la media dei valori nel suo vicinato rettangolare.  Il
kernel è composto da tutti 1, normalizzati per il numero di elementi (es.\ diviso 9
per un kernel $3\times3$).  Produce sfocatura (blurring).

\paragraph{Filtro Gaussiano.}
Un filtro passa-basso più sofisticato che usa un kernel con pesi estratti da una
distribuzione Gaussiana 2D.  I pixel più vicini al centro ricevono pesi maggiori:
\[
  G_\sigma(x,y) = \frac{1}{2\pi\sigma^2}\,e^{-\frac{x^2+y^2}{2\sigma^2}}
\]
Il parametro $\sigma$ controlla la quantità di sfocatura.

\begin{propbox}[Proprietà del filtro Gaussiano]
\begin{itemize}
  \item Convolvere un Gaussiano con un altro Gaussiano produce un nuovo Gaussiano più
        largo.
  \item Un kernel Gaussiano 2D è \textbf{separabile}: la convoluzione 2D può essere
        eseguita come due convoluzioni 1D (una orizzontale, una verticale), molto più
        efficienti computazionalmente.
\end{itemize}
\end{propbox}

\begin{notebox}[Costo computazionale dei filtri separabili]
Un filtro 2D è separabile quando la sua matrice kernel ha rango 1.  Per
un'immagine $M\times N$ e un kernel $L\times L$:
\begin{itemize}
  \item Convoluzione 2D standard: $\mathcal{O}(L^2 MN)$
  \item Implementazione separabile (due conv.\ 1D): $\mathcal{O}(2LMN)$
\end{itemize}
\end{notebox}

\paragraph{Filtro di nitidezza (unsharp masking).}
Un effetto di nitidezza si ottiene amplificando i dettagli dell'immagine.  I
dettagli vengono estratti sottraendo la versione sfocata all'originale, poi
riaggiunti:
\[
  F_{\text{sharp}} = F_{\text{orig}} + \alpha\bigl(F_{\text{orig}} - F_{\text{blur}}\bigr)
\]
dove $\alpha$ è la forza di nitidezza ($F_{\text{sharp}} = F_{\text{orig}}$ per
$\alpha=0$).  Si può implementare con un unico filtro che approssima un'impulso
scalato meno un Gaussiano.

% =============================================================================
\section{Filtraggio non lineare}
% =============================================================================

I filtri non lineari eseguono operazioni di vicinato, ma il nuovo valore del pixel
\emph{non} è una combinazione lineare dei suoi vicini.

% -- Bilateral ----------------------------------------------------------------
\subsection{Filtro bilaterale}

È un filtro di smoothing avanzato che \textbf{preserva i bordi}.  Un filtro
Gaussiano standard media i pixel indipendentemente dal loro contenuto, causando
sfocatura attraverso i bordi degli oggetti.  Il filtro bilaterale risolve questo
problema pesando i pixel vicini in base sia alla loro \textbf{distanza spaziale}
che alla loro \textbf{differenza di intensità}:
\[
  h[m,n] = \frac{1}{W_{mn}} \sum_{k,l}
    \underbrace{g[k,l]}_{\text{kernel spaziale}}\;
    \underbrace{r_{mn}[k,l]}_{\text{kernel di range}}\;
    f[m+k,\,n+l]
\]
dove:
\begin{itemize}
  \item $g[k,l]$: kernel Gaussiano spaziale, pesa i pixel in base alla distanza
        Euclidea dal pixel centrale $(m,n)$.
  \item $r_{mn}[k,l]$: kernel di range (o di intensità), pesa i pixel in base alla
        differenza di intensità dal pixel centrale $f[m,n]$.
  \item $W_{mn} = \sum_{k,l} g[k,l]\,r_{mn}[k,l]$: fattore di normalizzazione.
\end{itemize}

Un pixel contribuisce alla media solo se è sia spazialmente vicino sia
fotometricamente simile.  Il filtro bilaterale elimina il rumore nelle regioni piatte
ma preserva i bordi netti.  Applicazioni: denoising di alta qualità, ritocco
fotografico, effetto cartoon.

% -- Mediana ------------------------------------------------------------------
\subsection{Filtro mediana}

Sostituisce ogni pixel con il valore \textbf{mediano} di tutti i pixel nel suo
vicinato.

\begin{propbox}[Mediana vs.\ media]
Il filtro mediana rimuove il \textbf{rumore sale e pepe}: i valori anomali
(\emph{outlier}) vengono ignorati dalla mediana, mentre la media diffonde il rumore
nei pixel circostanti e sfoca i bordi.
\end{propbox}

% -- Sogliatura ---------------------------------------------------------------
\subsection{Sogliatura (Thresholding)}

Operazione non lineare semplice che converte un'immagine in scala di grigi in
un'immagine binaria:
\[
  g(m,n) = \begin{cases} 255 & \text{se } f(m,n) > A \\ 0 & \text{altrimenti} \end{cases}
\]

% =============================================================================
\section{Scalatura dell'immagine}
% =============================================================================

\subsection{Sottocampionamento e aliasing}

Per ridurre le dimensioni di un'immagine, un approccio ingenuo è il
\textbf{sottocampionamento}: eliminare righe e colonne.

\begin{warnbox}[Aliasing]
Il semplice sottocampionamento può causare \textbf{aliasing}: un effetto in cui
dettagli fini e pattern nell'immagine originale vengono distorte in nuovi pattern
errati nell'immagine ridotta.  Questo accade quando la \textbf{frequenza di
campionamento è troppo bassa} per catturare le frequenze più alte nell'immagine.

Per evitarlo, la frequenza di campionamento deve essere almeno il doppio della
frequenza massima presente nell'immagine (\textbf{frequenza di Nyquist}):
\[
  f_s \geq 2\,f_{\max}
\]
\end{warnbox}

\textbf{Anti-aliasing:} il modo corretto per ridurre un'immagine è applicare prima
un \textbf{filtro passa-basso} (es.\ blur Gaussiano) per rimuovere i dettagli ad
alta frequenza che non possono essere rappresentati alla risoluzione inferiore.  Dopo
la sfocatura, l'immagine può essere sottocampionata in modo sicuro.

\begin{notebox}[Aliasing nei sistemi pratici]
L'aliasing appare anche nelle implementazioni pratiche del downsampling nei framework
di deep learning.  Le operazioni di pooling nelle CNN possono ridurre
l'\emph{equivarianza} rispetto alle traslazioni, rendendo le predizioni instabili a
piccole traslazioni dell'immagine.  Il \emph{anti-aliased pooling} risolve il
problema inserendo un filtro passa-basso prima del sottocampionamento.
\end{notebox}

\subsection{Upsampling e interpolazione}

Per aumentare le dimensioni di un'immagine è necessario \textbf{interpolare} nuovi
pixel:
\begin{itemize}
  \item \textbf{Nearest-neighbor}: replica i pixel; il più semplice ma produce
        risultati a blocchi.
  \item \textbf{Bilineare}: calcola ciascun nuovo pixel come media pesata dei 4
        pixel più vicini nell'immagine originale; risultati più fluidi.
  \item \textbf{Bicubica}: usa un vicinato di 16 pixel per un risultato ancora più
        fluido.
\end{itemize}

\begin{notebox}[Super-risoluzione moderna]
Tecniche avanzate usano reti neurali per imparare come generare dettagli ad alta
risoluzione, oppure metodi multi-immagine che catturano foto leggermente shiftate e
le combinano per ricostruire un'immagine a risoluzione maggiore.
\end{notebox}

% =============================================================================
\section{Derivate dell'immagine e rilevamento dei bordi}
% =============================================================================

\begin{defbox}[Bordi]
I \textbf{bordi} (\emph{edges}) sono posizioni in un'immagine con un
\textbf{rapido cambiamento di intensità}.  Sono caratteristiche fondamentali per
comprendere il contenuto dell'immagine.
\end{defbox}

\subsection{Origine fisica dei bordi}

I bordi possono essere causati da diversi fenomeni fisici:
\begin{itemize}
  \item Discontinuità della normale alla superficie.
  \item Discontinuità di profondità.
  \item Discontinuità di colore o materiale della superficie.
  \item Discontinuità di illuminazione (es.\ ombre).
\end{itemize}

\subsection{Il gradiente dell'immagine}

Poiché un'immagine è una funzione 2D $f(x,y)$, possiamo analizzare le variazioni
di intensità calcolando le sue derivate, approssimate tramite convoluzione con
piccoli kernel.

\begin{itemize}
  \item La \textbf{derivata parziale} $\dfrac{\partial f}{\partial x}$ si approssima
        convolvendo con un filtro come $[-1,\,1]$ oppure $[1,\,-1]$.
        (Analogamente per $\dfrac{\partial f}{\partial y}$ in verticale.)

  \item Il \textbf{gradiente} $\nabla f = \left[\dfrac{\partial f}{\partial x},\,
        \dfrac{\partial f}{\partial y}\right]$ è un vettore che punta nella direzione
        del più rapido aumento di intensità.

  \item La \textbf{magnitudine del gradiente}:
        \[
          \norm{\nabla f} = \sqrt{\left(\frac{\partial f}{\partial x}\right)^2 +
                                  \left(\frac{\partial f}{\partial y}\right)^2}
        \]
        misura la forza del bordo.

  \item L'\textbf{orientazione del gradiente}:
        \[
          \theta = \mathrm{atan2}\!\left(\frac{\partial f}{\partial y},\,
                                         \frac{\partial f}{\partial x}\right)
        \]
        fornisce la direzione del bordo.
\end{itemize}

\textbf{Gestione del rumore.}  Applicare direttamente filtri derivativi è
altamente sensibile al rumore.  La soluzione è \textbf{prima uniformare l'immagine
con un filtro Gaussiano}.  Grazie alla proprietà associativa della convoluzione:
\[
  \frac{\partial}{\partial x}(G * f)
  = \left(\frac{\partial}{\partial x}G\right) * f
\]
Si ottiene il filtro \textbf{Derivative of Gaussian (DoG)}: invece di sfocare
l'immagine e poi prenderne la derivata, si prende la derivata del kernel Gaussiano
e si convolve con essa l'immagine in un'unica operazione efficiente.

\begin{warnbox}[Nota terminologica: DoG]
In rilevamento dei bordi, l'espressione \emph{Derivative of Gaussian} non va
abbreviata in DoG, perché in rilevamento di feature e SIFT, DoG significa
tipicamente \emph{Difference of Gaussians}.  Si tratta di concetti distinti.
\end{warnbox}

\subsection{Derivate del secondo ordine: il Laplaciano}

La derivata del secondo ordine può anch'essa essere usata per trovare i bordi.
Il \textbf{Laplaciano} di un'immagine è:
\[
  \nabla^2 f = \frac{\partial^2 f}{\partial x^2} + \frac{\partial^2 f}{\partial y^2}
\]
che è uno scalare.

\begin{itemize}
  \item In un segnale 1D, un bordo corrisponde a un picco nella prima derivata e a
        un \textbf{zero-crossing} nella seconda derivata.
  \item Combinando il Laplaciano con lo smoothing Gaussiano si crea un filtro
        \textbf{Laplacian of Gaussian (LoG)}.  I bordi vengono localizzati trovando
        gli zero-crossing nell'immagine filtrata.
\end{itemize}

\begin{notebox}[Zero-crossing vs.\ DoG]
Gli zero-crossing (LoG) sono più accurati nel localizzare i bordi, ma il DoG è
computazionalmente più efficiente.  La prima derivata in corrispondenza di un bordo
corrisponde a un massimo, mentre la seconda derivata è zero.
\end{notebox}

\subsection{L'operatore di Sobel}

L'operatore di Sobel è un'approssimazione efficiente del DoG.  È un
\textbf{filtro separabile} composto da un kernel derivativo 1D e un kernel di
smoothing 1D.  Il filtro Sobel orizzontale combina $[1,\,2,\,1]$ (smoothing) e
$[1,\,0,\,-1]$ (derivativo).

Altri filtri derivativi comuni: \textbf{Prewitt}, \textbf{Roberts}, \textbf{Scharr}.

\begin{notebox}[Normalizzazione di Sobel]
I mask di Sobel standard sono spesso scritti senza il fattore di normalizzazione
$1/8$.  Questo non influisce sul rilevamento dei bordi se si usano solo risposte
relative o soglie, ma il fattore è necessario per ottenere la scala corretta della
magnitudine del gradiente.
\end{notebox}

\subsection{La pipeline di Canny}

Il rilevatore di Canny è un algoritmo multi-stadio considerato il metodo classico
standard per il rilevamento dei bordi.

\begin{enumerate}
  \item \textbf{Filtraggio}: si uniforma l'immagine con un filtro Gaussiano e si
        calcolano magnitudine e orientazione del gradiente usando un filtro DoG (o
        un'approssimazione come Sobel).

  \item \textbf{Non-maximum suppression}: per ogni pixel, si verifica se la sua
        magnitudine è un massimo locale nella direzione del suo gradiente; se no, lo
        si sopprime impostandolo a 0.  Ciò produce bordi sottili, larghi un pixel.

  \item \textbf{Linking e sogliatura isteresi}: si usano due soglie (alta e bassa).
        \begin{itemize}
          \item Pixel con magnitudine del gradiente sopra la soglia alta:
                bordi \emph{forti}, conservati.
          \item Pixel con magnitudine tra soglia bassa e alta: bordi \emph{deboli},
                conservati solo se connessi a un pixel di bordo forte.
          \item Pixel sotto la soglia bassa: scartati.
        \end{itemize}
\end{enumerate}

\begin{propbox}[Criteri di ottimalità del rilevatore di Canny]
Un rilevatore di bordi ottimale deve soddisfare tre criteri:
\begin{enumerate}
  \item \textbf{Buona rilevazione}: pochi falsi positivi e falsi negativi.
  \item \textbf{Buona localizzazione}: i bordi rilevati devono essere vicini alla
        vera posizione del bordo.
  \item \textbf{Risposta singola}: un bordo fisico deve generare un solo bordo
        rilevato.
\end{enumerate}
Il comportamento dipende principalmente da tre parametri: la scala di smoothing
Gaussiano $\sigma$, la soglia alta e la soglia bassa.  Un $\sigma$ più grande
sopprime i dettagli fini e rileva bordi su scala maggiore; un $\sigma$ più piccolo
preserva bordi più fini ma è più sensibile al rumore.
\end{propbox}

% =============================================================================
\chapter{Piramidi di immagini}
% =============================================================================

\section{Piramidi Gaussiane}

\begin{defbox}[Piramide Gaussiana]
Una \textbf{piramide Gaussiana} è una rappresentazione multi-risoluzione
dell'immagine in cui ogni livello è creato applicando un filtro Gaussiano e
sottocampionando il livello precedente di un fattore 2, producendo immagini
progressivamente più piccole e sfocate, utili per la \textbf{ricerca
coarse-to-fine}.
\end{defbox}

Man mano che si sale di livello nella piramide (risoluzioni minori), i dettagli ad
alta frequenza vengono eliminati, mentre le componenti a bassa frequenza e le regioni
uniformi vengono preservate.

\begin{notebox}[Perché si usa il filtro Gaussiano]
Il Gaussiano è usato perché è \textbf{auto-riproducente}: convolvere un Gaussiano
con un altro Gaussiano produce un nuovo Gaussiano con varianza maggiore.  Questo
rende possibile lo smoothing incrementale tra i livelli della piramide.  Una volta
filtrata passa-basso l'immagine, sono necessari meno pixel per rappresentare il
contenuto rimanente senza aliasing severo.
\end{notebox}

\begin{warnbox}[La sfocatura è con perdita di informazione]
Non è possibile ricostruire l'immagine originale ad alta risoluzione a partire da
un livello a risoluzione inferiore.
\end{warnbox}

\begin{notebox}[Dimensione di una piramide Gaussiana]
Se ogni livello è sottocampionato di un fattore 2, sia la larghezza che l'altezza
vengono dimezzate $\Rightarrow$ il livello successivo ha $1/4$ dei pixel del
precedente.  Il numero totale di pixel dell'intera piramide:
\[
  N_{\text{tot}} = N + \frac{N}{4} + \frac{N}{16} + \cdots
  = N\sum_{i=0}^{\infty}\left(\frac{1}{4}\right)^i
  = N \cdot \frac{1}{1-1/4}
  = \frac{4}{3}N
\]
L'intera piramide richiede solo circa il $\frac{4}{3}$ dello storage
dell'immagine originale.
\end{notebox}

\section{Piramidi Laplaciane}

\begin{defbox}[Piramide Laplaciana]
La \textbf{piramide Laplaciana} è una \textbf{rappresentazione senza perdita} che
memorizza i dettagli ad alta frequenza persi tra due livelli di una piramide
Gaussiana, calcolando un'immagine \emph{residuo}.
\end{defbox}

Si definiscono due operazioni fondamentali:

\begin{itemize}
  \item $\mathrm{REDUCE}(I)$: crea il livello successivo $G_{i+1}$ della piramide
        Gaussiana sfocando e poi sottocampionando $G_i$:
        \[
          G_{i+1} = \mathrm{REDUCE}(G_i)
          = \mathrm{downsample}\!\bigl(\mathrm{blur}(G_i)\bigr)
        \]

  \item $\mathrm{EXPAND}(I)$: predice un'immagine a risoluzione maggiore $G_i$ da
        $G_{i+1}$ a risoluzione minore, tramite upsampling e sfocatura:
        \[
          \mathrm{EXPAND}(G_{i+1})
          = \mathrm{blur}\!\bigl(\mathrm{upsample}(G_{i+1})\bigr)
        \]
\end{itemize}

Il \textbf{residuo} $L_i$ forma un livello della piramide Laplaciana:
\[
  L_i = G_i - \mathrm{EXPAND}(G_{i+1})
\]

Ogni residuo contiene principalmente i dettagli persi nel passaggio da un livello
Gaussiano a quello successivo (bordi, contorni, texture fine).

La rappresentazione completa della piramide Laplaciana è
$\{L_0,L_1,\ldots,L_{n-1}\}$ più il livello Gaussiano più piccolo $G_n$.

\textbf{Perché ``Laplaciana''?}  La creazione di ogni residuo (immagine meno versione
sfocata, differenza di Gaussiani) approssima l'operatore LoG:
\[
  \underbrace{G_i - \mathrm{EXPAND}(G_{i+1})}_{\text{Difference of Gaussians}}
  \approx
  \underbrace{\nabla^2 G}_{\text{Laplacian of Gaussian}}
\]

\subsection{Sommario delle formule}

\[
\boxed{
  G_{i+1} = \mathrm{REDUCE}(G_i) \qquad
  L_i = G_i - \mathrm{EXPAND}(G_{i+1}) \qquad
  G_i = L_i + \mathrm{EXPAND}(G_{i+1})
}
\]

\subsection{Ricostruzione}

L'immagine originale può essere \textbf{ricostruita perfettamente} dalla piramide
Laplaciana invertendo il processo:
\begin{enumerate}
  \item Si parte dall'immagine al livello più alto, $G_n$.
  \item Per ogni livello $i$ da $n-1$ fino a 0:
        \[G_i = \mathrm{EXPAND}(G_{i+1}) + L_i\]
  \item Il risultato finale $G_0$ è l'immagine originale perfettamente ricostruita.
\end{enumerate}

In assenza di errori di quantizzazione e approssimazione numerica, la piramide
Laplaciana consente \textbf{ricostruzione perfetta}.  Per memorizzare questa
ricostruzione si devono conservare: le immagini residuo a ogni livello; l'immagine
più piccola al vertice della piramide Gaussiana.

\section{Applicazioni delle piramidi}

Le piramidi di immagini sono uno strumento fondamentale usato in: compressione di
immagini, denoising, registrazione multi-scala (allineamento immagini da grossolano
a fine) e rilevamento multi-scala (trovare oggetti a varie scale).

\subsection{Blending di immagini}

Le piramidi Laplaciane forniscono una tecnica potente per fondere due immagini in
modo fluido, evitando la giunzione netta che risulterebbe da un semplice taglia-e-incolla.

\begin{notebox}[Algoritmo di blending piramidale]
Per fondere due immagini $A$ e $B$ usando un'immagine maschera $M$:
\begin{enumerate}
  \item Costruire le piramidi Laplaciane per le due immagini sorgente, $LA$ e $LB$.
  \item Costruire una piramide Gaussiana per la maschera, $GM$.  I valori della
        maschera (da 0 a 1) fungono da pesi di blending a ogni scala.
  \item Per ogni livello $j$, creare un livello Laplaciano combinato $LC_j$ come
        media pesata:
        \[
          LC_j(x,y)
          = GM_j(x,y)\cdot LA_j(x,y)
          + \bigl(1-GM_j(x,y)\bigr)\cdot LB_j(x,y)
        \]
  \item Collassare la piramide Laplaciana combinata $LC$ per ottenere l'immagine
        finale.
\end{enumerate}
\end{notebox}

\subsection{Altre rappresentazioni multi-risoluzione}

Altre rappresentazioni multi-risoluzione includono le \textbf{piramidi orientabili}
(\emph{steerable pyramids}) e le \textbf{wavelet}.
\begin{itemize}
  \item Le piramidi orientabili memorizzano risposte orientate multiple a ogni scala.
  \item Le wavelet decompongono l'immagine in approssimazioni a bassa frequenza e
        dettagli ad alta frequenza, spesso separati in componenti orizzontali,
        verticali e diagonali.
\end{itemize}

\subsection{Rilevamento multi-scala}

Nel rilevamento multi-scala, lo stesso template o rilevatore a dimensione fissa
viene applicato a livelli diversi della piramide.  Poiché gli oggetti appaiono a
dimensioni diverse a seconda del livello, un rilevatore che funziona a una dimensione
fissa di template può rilevare oggetti di diverse dimensioni apparenti scansionando
la piramide.

% =============================================================================
\chapter{Dominio delle frequenze}
% =============================================================================

\section{Dominio delle frequenze e serie di Fourier}

Qualsiasi funzione univariata può essere riscritta come somma pesata di seni e coseni
di frequenze diverse, sotto certe condizioni.

\subsection{Costruire segnali da sinusoidi}

Il mattone di base per l'analisi di Fourier è la sinusoide:
\[
  A\sin(\omega x + \phi)
\]
definita da tre parametri chiave:
\begin{itemize}
  \item \textbf{Ampiezza} ($A$): la forza o intensità dell'onda.
  \item \textbf{Frequenza angolare} ($\omega$): determina la velocità di oscillazione.
        Il periodo è $T = 2\pi/|\omega|$.
  \item \textbf{Fase} ($\phi$): la posizione iniziale o offset della sinusoide.
\end{itemize}

Il \textbf{teorema di Fourier} afferma che sommando abbastanza sinusoidi con
frequenze e ampiezze diverse, si può approssimare qualsiasi segnale periodico.

Lo \textbf{spettro di frequenza} è un grafico che mostra l'ampiezza di ciascuna
componente di frequenza presente nel segnale.  Il valore alla frequenza zero è detto
\textbf{componente DC} e rappresenta il valore medio del segnale.

\section{La trasformata di Fourier}

Le serie di Fourier rappresentano funzioni periodiche.  La \textit{trasformata di
Fourier} estende questa idea alle \textbf{funzioni non periodiche}:

\begin{defbox}[Trasformata di Fourier]
La trasformata di Fourier converte un segnale $f(x)$ dal dominio spaziale a una
funzione $F(\omega)$ nel dominio delle frequenze.  $F(\omega)$ è a valori complessi:
\[
  F(\omega) = R(\omega) + i\,I(\omega)
\]
L'ampiezza e la fase si ricavano dalla parte reale e immaginaria:
\[
  A(\omega) = \sqrt{R(\omega)^2 + I(\omega)^2}, \qquad
  \phi(\omega) = \mathrm{atan2}\!\bigl(I(\omega),\,R(\omega)\bigr)
\]
L'uso di $\mathrm{atan2}$ è preferibile all'arcotangente semplice perché gestisce
correttamente il quadrante del numero complesso.  La trasformazione è invertibile.
\end{defbox}

\subsection{Definizione matematica}

\textbf{Trasformata diretta} (Spaziale $\to$ Frequenza):
\[
  F(\omega) = \int_{-\infty}^{+\infty} f(x)\,e^{-i\omega x}\,dx
\]

\textbf{Trasformata inversa} (Frequenza $\to$ Spaziale):
\[
  f(x) = \frac{1}{2\pi}\int_{-\infty}^{+\infty} F(\omega)\,e^{i\omega x}\,d\omega
\]

La connessione con le sinusoidi è data dalla \textbf{formula di Eulero}:
$e^{i\theta} = \cos\theta + i\sin\theta$.
La trasformata misura essenzialmente quanta di una data frequenza seno e coseno è
presente nel segnale $f(x)$.

\subsection{La trasformata discreta di Fourier (DFT)}

Per segnali e immagini digitali si usa la \textit{Discrete Fourier Transform (DFT)}.
Per un segnale 1D $f(x)$ di lunghezza $N$:
\[
  F(k) = \sum_{x=0}^{N-1} f(x)\,e^{-j2\pi kx/N}
\]
Può essere espressa come moltiplicazione matriciale $\mathbf{F} = \mathbf{W}\mathbf{f}$.

\begin{notebox}[Fast Fourier Transform (FFT)]
Il calcolo diretto della DFT tramite moltiplicazione matriciale ha complessità
$\mathcal{O}(N^2)$.  La \textbf{Fast Fourier Transform (FFT)} calcola la DFT in
$\mathcal{O}(N\log N)$, rendendola pratica per l'elaborazione delle immagini.
\end{notebox}

\section{Analisi di Fourier in 2D}

La trasformata di Fourier si estende a 2D per analizzare le immagini.  Un'immagine
$f(x,y)$ viene trasformata nella sua rappresentazione in frequenza $F(u,v)$, dove
$u$ e $v$ sono le frequenze orizzontale e verticale:
\[
  F(u,v) = \sum_{x=0}^{M-1}\sum_{y=0}^{N-1}
    f(x,y)\,e^{-j2\pi\!\left(\frac{ux}{M}+\frac{vy}{N}\right)}
\]

Nello spettro 2D di Fourier (visualizzato con il DC al centro):
\begin{itemize}
  \item Le \textbf{basse frequenze} sono vicino al centro e corrispondono a strutture
        lisce su larga scala.
  \item Le \textbf{alte frequenze} sono lontane dal centro e corrispondono a dettagli
        fini, bordi e texture.
  \item L'orientazione delle linee nello spettro corrisponde all'orientazione dei
        bordi nell'immagine (es.\ i bordi verticali creano linee orizzontali nel
        dominio delle frequenze).
\end{itemize}

\begin{notebox}[L'importanza della fase]
\begin{itemize}
  \item L'\textbf{ampiezza} della DFT rappresenta quanta di una particolare frequenza
        è presente nell'immagine $\Rightarrow$ manca di informazione spaziale.
  \item La \textbf{fase} della DFT rappresenta la disposizione spaziale di ogni
        componente di frequenza $\Rightarrow$ definisce la posizione di bordi, linee,
        angoli e texture.
  \item Se si ricostruisce un'immagine usando solo la fase, il risultato è
        riconoscibile ma distorto, con strane variazioni di intensità.  La fase è
        generalmente più importante dell'ampiezza per la ricostruzione visiva.
\end{itemize}
\end{notebox}

\section{Coppie della trasformata di Fourier e proprietà}

\textbf{Coppie importanti:}
\begin{itemize}
  \item Funzione \emph{box} nel dominio spaziale $\leftrightarrow$ funzione
        \emph{sinc} nel dominio delle frequenze.
  \item Gaussiano nel dominio spaziale $\leftrightarrow$ Gaussiano nel dominio delle
        frequenze.
  \item Per dualità: sinc nel dominio spaziale $\leftrightarrow$ box nel dominio
        delle frequenze.
\end{itemize}

\textbf{Proprietà di scaling:} espandere un segnale nel dominio spaziale comprime
la sua trasformata nel dominio delle frequenze, e viceversa.  Strutture lisce su
larga scala $\leftrightarrow$ basse frequenze; dettagli piccoli e variazioni rapide
$\leftrightarrow$ alte frequenze.

\section{Filtraggio nel dominio delle frequenze}

\subsection{Il teorema della convoluzione}

\begin{propbox}[Teorema della convoluzione]
La convoluzione nel dominio spaziale è equivalente alla moltiplicazione
element-wise nel dominio delle frequenze:
\[
  \mathcal{F}\{g * h\} = \mathcal{F}\{g\} \cdot \mathcal{F}\{h\}
\]
Equivalentemente, la moltiplicazione nel dominio delle frequenze corrisponde alla
convoluzione nel dominio spaziale.
\end{propbox}

Invece di eseguire una costosa convoluzione nel dominio spaziale, si può:
\begin{enumerate}
  \item Trasformare immagine e kernel nel dominio delle frequenze tramite FFT.
  \item Moltiplicarli element-wise.
  \item Ritrasformare il risultato nel dominio spaziale tramite FFT inversa.
\end{enumerate}

\textbf{Esempi di filtraggio:}
\begin{itemize}
  \item \textbf{Filtraggio passa-basso}: moltiplicare lo spettro per una maschera
        che preserva le basse frequenze (vicino al centro) e attenua le alte
        frequenze.  Un Gaussiano è un eccellente filtro passa-basso.
  \item \textbf{Filtraggio passa-alto}: attenuare le basse frequenze e preservare le
        alte (per nitidezza o rilevamento bordi).
  \item \textbf{Filtraggio passa-banda}: preservare un anello di frequenze per
        isolare feature a una scala specifica.
\end{itemize}

\begin{notebox}[Filtro box vs.\ filtro Gaussiano]
Il \textbf{filtro box} è brusco nel dominio spaziale.  La sua trasformata di
Fourier è una \textbf{funzione sinc} con sidelobes significativi.  Queste
caratteristiche non ideali producono un taglio di frequenza imperfetto e possono
causare \textbf{artefatti di ringing} nell'immagine sfocata risultante.

Il \textbf{filtro Gaussiano} è regolare sia nel dominio spaziale che in quello
delle frequenze, garantendo un'attenuazione pulita.  Un filtro passa-basso ideale
(taglio netto in frequenza) produce anch'esso ringing perché la sua risposta
impulsiva nel dominio spaziale è una sinc.

In pratica, il supporto spaziale di un mask Gaussiano viene scelto pari a circa
\textbf{tre volte la deviazione standard}, in modo da includere la maggior parte
della campana Gaussiana.
\end{notebox}

\section{Campionamento}

\begin{propbox}[Teorema di Nyquist-Shannon]
Il teorema di Nyquist-Shannon afferma che un segnale continuo può essere
ricostruito perfettamente dai suoi campioni discreti \textbf{senza aliasing} se e
solo se:
\[
  f_s \geq 2\,f_{\max}
\]
dove la frequenza minima di campionamento $2f_{\max}$ è chiamata
\textbf{frequenza di Nyquist}.
\end{propbox}

\begin{notebox}[Applicazione alle piramidi Gaussiane]
La sfocatura Gaussiana (filtraggio passa-basso) riduce $f_{\max}$, prevenendo
così l'aliasing quando l'immagine viene successivamente sottocampionata.
\end{notebox}

\section{Immagini ibride}

Le \textbf{immagini ibride} sono create combinando le componenti a bassa frequenza
di un'immagine con le componenti ad alta frequenza di un'altra.

\begin{enumerate}
  \item Un \textbf{filtro passa-basso} (es.\ Gaussiano) viene applicato alla prima
        immagine, conservandone solo la struttura grossolana.
  \item Un \textbf{filtro passa-alto} (es.\ Laplaciano) viene applicato alla seconda
        immagine, conservandone solo i dettagli fini e i bordi.
  \item Le due immagini risultanti vengono sommate.
\end{enumerate}

L'immagine ibrida finale viene percepita diversamente a seconda della distanza di
visione: da lontano il sistema visivo coglie l'immagine a bassa frequenza; da vicino
i dettagli ad alta frequenza diventano dominanti.  Questo sfrutta il comportamento
\textbf{dipendente dalla frequenza della visione umana}.

% =============================================================================
\chapter{Feature locali}
% =============================================================================

\section{Introduzione alle feature locali}

\begin{defbox}[Feature]
Una \textbf{feature} è un'informazione sul contenuto di un'immagine.  Si distinguono
due tipi principali:
\begin{itemize}
  \item \textbf{Feature globali}: descrivono un'immagine nel suo insieme per
        rappresentare un intero oggetto.  Usate per il rilevamento di oggetti
        (l'oggetto è presente, sì/no?).
  \item \textbf{Feature locali}: descrivono patch o keypoint all'interno di un
        oggetto.  Usate per il riconoscimento di oggetti, dove è necessario trovare
        corrispondenze tra parti degli oggetti.
\end{itemize}
\end{defbox}

\begin{propbox}[Proprietà di buone feature locali]
Buone feature locali devono essere \textbf{numerose} per il matching e
\textbf{distintive} per una corrispondenza non ambigua.  Crucialmente, devono essere
\textbf{invarianti} alle trasformazioni fotometriche e geometriche.

Le feature locali sono utili perché sono: \textbf{locali} (robuste a occlusione
parziale e cluttering), \textbf{numerose}, \textbf{distintive} (permettono un
matching non ambiguo) ed \textbf{efficienti} (adatte ad applicazioni in tempo reale).
\end{propbox}

\section{Invarianza ed equivarianza}

\begin{itemize}
  \item \textbf{Invarianza}: una trasformazione dell'immagine non cambia la proprietà
        della feature rilevata.
  \item \textbf{Equivarianza}: se l'immagine è trasformata geometricamente, le
        posizioni delle feature rilevate si trasformano di conseguenza, così i punti
        corrispondenti vengono ancora rilevati in posizioni corrispondenti.
\end{itemize}

Per i rilevatori di feature locali, si vuole tipicamente \textbf{invarianza} alle
trasformazioni fotometriche ed \textbf{equivarianza} alle trasformazioni
geometriche.

\section{Applicazioni delle feature locali}

Feature locali sono usate in: allineamento di immagini e stitching di panorami,
ricostruzione 3D, SLAM visivo, tracking in realtà aumentata, riconoscimento di
oggetti, image retrieval e navigazione robotica/veicolare.

\section{Detection, description e matching}

I metodi a feature locali sono tipicamente composti da tre fasi:
\begin{itemize}
  \item \textbf{Detection}: identificare punti di interesse o keypoint nell'immagine.
  \item \textbf{Description}: estrarre un vettore descrittore dalla patch intorno a
        ogni keypoint.
  \item \textbf{Matching}: confrontare i descrittori tra immagini diverse per trovare
        corrispondenze.
\end{itemize}

\section{Rilevatore di angoli di Harris}

\begin{defbox}[Harris corner detector]
Il \textit{rilevatore di angoli di Harris} identifica punti importanti traslando una
piccola finestra intorno a essi e verificando dove si verificano forti cambiamenti
di aspetto.
\end{defbox}

\subsection{La matematica del rilevamento di angoli}

Per quantificare il cambiamento di intensità, si calcola la \textbf{Sum of Squared
Differences (SSD)} per una finestra $W$ traslata di un piccolo offset $(u,v)$:
\[
  E(u,v) = \sum_{(x,y)\in W} \bigl[I(x+u,\,y+v) - I(x,y)\bigr]^2
\]

Un angolo è un punto in cui questo errore $E(u,v)$ è elevato per tutti i piccoli
spostamenti $(u,v)$.

\textbf{Approssimazione per piccoli spostamenti.}  Per piccoli spostamenti si
approssima il cambiamento di intensità con un'espansione di Taylor al primo ordine:
\[
  I(x+u,\,y+v) \approx I(x,y) + I_x\,u + I_y\,v
\]
Sostituendo nella formula della SSD:
\[
  E(u,v) \approx \sum_{(x,y)\in W}\bigl[I_x\,u + I_y\,v\bigr]^2
  \approx Au^2 + 2Buv + Cv^2
\]
dove $A = \sum I_x^2$, $B = \sum I_x I_y$, $C = \sum I_y^2$
(somme sulla finestra $W$).

\subsection{La matrice dei secondi momenti}

La forma quadratica si esprime in notazione matriciale:
\[
  E(u,v) \approx
  \begin{bmatrix}u & v\end{bmatrix}
  \underbrace{\begin{bmatrix}A & B \\ B & C\end{bmatrix}}_{\mathbf{H}}
  \begin{bmatrix}u \\ v\end{bmatrix}
\]

La matrice $\mathbf{H}$ è la \textbf{matrice dei secondi momenti} (o
\emph{structure tensor}):
\[
  \mathbf{H} = \sum_{(x,y)\in W}
  \begin{bmatrix}I_x^2 & I_x I_y \\ I_x I_y & I_y^2\end{bmatrix}
\]

\subsection{Interpretazione della matrice dei secondi momenti}

La forma della superficie di errore quadratica $E(u,v)$ è determinata dagli
autovalori $\lambda_1$ e $\lambda_2$ di $\mathbf{H}$:

\begin{itemize}
  \item Gli \textbf{autovettori} di $\mathbf{H}$ corrispondono alle direzioni di
        massimo e minimo cambiamento di intensità.
  \item Gli \textbf{autovalori} $\lambda_{\max}$ e $\lambda_{\min}$ rappresentano
        la magnitudine di tale cambiamento in quelle direzioni.
\end{itemize}

Questo permette di classificare la regione:
\begin{itemize}
  \item \textbf{Flat}: entrambi $\lambda_1$ e $\lambda_2$ piccoli.  Errore basso in
        tutte le direzioni.
  \item \textbf{Edge}: un autovalore grande e l'altro piccolo
        ($\lambda_1 \gg \lambda_2$).  Errore alto in una direzione ma basso lungo
        il bordo.
  \item \textbf{Corner}: entrambi $\lambda_1$ e $\lambda_2$ grandi e di magnitudine
        simile.  Errore alto in tutte le direzioni.
\end{itemize}

\begin{notebox}[Criterio del minimo autovalore]
Una misura semplice di \emph{cornerness} è $\lambda_{\min}$, il minore autovalore
della matrice dei secondi momenti.  Un punto è considerato un buon angolo se
$\lambda_{\min}$ supera una soglia.
\end{notebox}

\subsection{L'operatore di Harris e l'algoritmo}

Calcolare gli autovalori per ogni pixel è inefficiente.  Il rilevatore di Harris usa
una \textbf{funzione di risposta} $R$ che approssima il comportamento degli autovalori
senza calcolarli esplicitamente:
\[
  R = \det(\mathbf{H}) - k\cdot\trace^2(\mathbf{H})
  = \lambda_1\lambda_2 - k(\lambda_1+\lambda_2)^2
\]

dove $k$ è una costante empirica (tipicamente 0.04--0.06).  Il valore di $R$ indica
la \emph{cornerness}: grande e positivo per un angolo, grande e negativo per un
bordo, piccolo per una regione piatta.

\textbf{Algoritmo del rilevatore di Harris:}
\begin{enumerate}
  \item Calcolare le derivate dell'immagine $I_x$ e $I_y$ (es.\ usando filtri
        Sobel).
  \item Per ogni pixel, calcolare la matrice dei secondi momenti $\mathbf{H}$ su una
        finestra Gaussiana intorno ad esso.
  \item Calcolare la risposta di Harris $R$ per ogni pixel.
  \item Sogliare la mappa di risposta $R$ per identificare candidati angoli forti.
  \item Eseguire la \textbf{non-maximum suppression} per trovare i punti angolo
        finali.
\end{enumerate}

\subsection{Proprietà e limitazioni}

Il rilevatore di Harris è:
\begin{itemize}
  \item \textbf{Equivariante} alla rotazione (l'ellisse dei secondi momenti ruota,
        ma gli autovalori sono preservati) e alla traslazione.
  \item \textbf{Invariante} agli shift di luminosità ($+b$) ma solo parzialmente ai
        cambiamenti di contrasto ($\times a$), poiché la sua risposta $R$ scala,
        influenzando la sogliatura.
  \item \textbf{Non invariante alla scala}: un angolo può apparire piatto o come un
        bordo quando si esegue lo zoom.
\end{itemize}

\section{Rilevamento di blob invariante alla scala}

Per superare la limitazione di scala di Harris, i \textit{blob} invarianti alla
scala vengono rilevati trovando la loro \textbf{scala caratteristica}: la scala
$\sigma$ in uno \textbf{scale-space 3D} $(x, y, \sigma)$ in cui una funzione di
risposta è massimizzata.

Il \textbf{LoG normalizzato alla scala} è un comune rilevatore di blob:
\[
  \nabla^2_{\text{norm}}\,g
  = \sigma^2\left(\frac{\partial^2 g}{\partial x^2}
                 +\frac{\partial^2 g}{\partial y^2}\right)
\]

Massimizzare questa risposta LoG attraverso le scale fornisce la scala caratteristica.

\begin{defbox}[Scale-space extrema]
Il rilevamento invariante alla scala può essere formulato come ricerca dei massimi
locali in uno scale-space 3D $(x, y, \sigma)$.  Un keypoint viene selezionato non
solo se la sua risposta è massima localmente nel piano immagine, ma anche se è
massima tra le scale vicine.  Il fattore $\sigma^2$ nella LoG normalizzata è
necessario per rendere le risposte comparabili tra scale diverse.
\end{defbox}

\begin{notebox}[Implementazione efficiente]
Invece di applicare filtri sempre più grandi all'immagine, è più efficiente
costruire una piramide Gaussiana e applicare un filtro di dimensione fissa a ogni
livello:
\begin{enumerate}
  \item Costruire uno scale-space calcolando la risposta LoG per un range di scale
        $\sigma$.
  \item Trovare i massimi locali nello spazio 3D posizione-scala $(x,y,\sigma)$.
  \item La posizione $(x,y)$ e la scala caratteristica $\sigma$ di ogni massimo
        definiscono una feature rilevata.
\end{enumerate}
\end{notebox}

\begin{notebox}[Filtro LoG a dimensione fissa e scala effettiva]
Quando si usa una piramide Gaussiana per il rilevamento di blob invariante alla
scala con un filtro LoG a dimensione fissa e $\sigma$ interno fisso:
\begin{itemize}
  \item \textbf{Livello 0} (immagine originale): scala effettiva
        $\approx \sigma_{\text{filter}}$.
  \item \textbf{Livello 1} (sottocampionato $2\times$): ogni pixel del filtro copre
        $2\times2$ pixel dell'immagine originale; scala effettiva
        $\approx 2\times\sigma_{\text{filter}}$.
  \item \textbf{Livello 2} (sottocampionato $4\times$): scala effettiva
        $\approx 4\times\sigma_{\text{filter}}$.
\end{itemize}
La scala caratteristica per un blob rilevato è la scala del livello piramidale in
cui la sua risposta LoG è massimizzata.
\end{notebox}

% =============================================================================
\chapter{Descrittori locali}
% =============================================================================

\section{Progettare descrittori di feature}

\begin{defbox}[Descrittore di feature]
Un \textbf{descrittore di feature} è una \textbf{rappresentazione vettoriale} della
patch che circonda un keypoint. L'obiettivo è progettare descrittori tali che
\textbf{patch con contenuto simile abbiano descrittori simili}, abilitando un
matching robusto. Un buon descrittore deve essere robusto alle trasformazioni
fotometriche e geometriche: il design è un compromesso tra \textbf{distintività}
e \textbf{invarianza}.
\end{defbox}

\subsection{Evoluzione di un descrittore}

\begin{enumerate}
  \item \textbf{Patch grezza}: intensità dei pixel; perfetta per match esatti ma
        altamente sensibile a variazioni di luminosità e posizione.
  \item \textbf{Gradienti dell'immagine}: invarianti agli shift di luminosità ma
        sensibili a deformazione, rotazione e cambiamenti di contrasto.
  \item \textbf{Istogramma colore}: invariante a rotazione/scala ma perde tutto
        il layout spaziale, rendendolo non discriminativo.
  \item \textbf{Istogrammi spaziali}: istogrammi per cella (come in HOG o nelle
        regioni di SIFT). Parzialmente invarianti alla deformazione; sensibili alla
        rotazione a meno che non si stabilisca prima un'orientazione dominante.
  \item \textbf{Normalizzazione dell'orientazione}: raggiunge la piena invarianza
        alla rotazione allineando la patch alla sua orientazione dominante prima
        del calcolo del descrittore (come in SIFT).
\end{enumerate}

\subsection{Confronto di istogrammi}

\begin{itemize}
  \item \textbf{Intersezione}: $\cap(Q,V)=\sum_i \min(q_i,v_i)$. Robusta.
  \item \textbf{Distanza Euclidea}: $d(Q,V)=\sqrt{\sum_i(q_i-v_i)^2}$.
        Meno robusta: pesa tutti i bin ugualmente.
  \item \textbf{Chi-quadro}: $\chi^2(Q,V)=\sum_i\frac{(q_i-v_i)^2}{q_i+v_i}$.
        Più discriminativa: pesa i bin diversamente.
\end{itemize}

\section{Descrittore HOG (Histogram of Oriented Gradients)}

\begin{defbox}[HOG]
\textbf{HOG} è un potente descrittore di feature usato per il rilevamento di
pedoni, che combina gradiente dell'immagine, celle spaziali e normalizzazione
locale del contrasto.
\end{defbox}

Il descrittore viene calcolato su una patch di dimensione fissa
(es.\ $64\times128$ pixel per il rilevamento di pedoni).

\paragraph{1.\ Calcolo del gradiente.}
Si calcolano i gradienti orizzontale e verticale per ottenere magnitudine e
direzione per ogni pixel, usando \textbf{gradienti unsigned} con angoli da $0^{\circ}$
a $180^{\circ}$ (direzioni opposte trattate ugualmente).

\paragraph{2.\ Istogrammi per cella.}
La patch viene divisa in piccole regioni dette \textbf{celle} (es.\ $8\times8$
pixel). Per ogni cella si crea un istogramma delle orientazioni del gradiente
con \textbf{9 bin} (per angoli $0^{\circ}, 20^{\circ}, \ldots, 160^{\circ}$). Ogni pixel vota con
peso uguale alla sua magnitudine, nel bin corrispondente alla sua orientazione.
I voti vengono interpolati bilinearmente tra bin adiacenti per evitare aliasing.

\paragraph{3.\ Normalizzazione per blocco.}
Le celle adiacenti vengono raggruppate in \textbf{blocchi} (es.\ $2\times2$ celle,
$16\times16$ pixel). Gli istogrammi di tutte le celle nel blocco vengono
concatenati e normalizzati:
\[
  v_{\mathrm{norm}} = \frac{v}{\sqrt{\|v\|_2^2+\varepsilon}}
\]
Questo rende il descrittore \textbf{robusto ai cambiamenti fotometrici locali}.

\paragraph{4.\ Vettore di feature.}
Tutti i vettori di blocco normalizzati vengono concatenati nel descrittore HOG
finale. I blocchi si sovrappongono (finestra mossa di 8 pixel alla volta),
fornendo una rappresentazione ridondante e robusta. Le celle d'angolo vengono
codificate una volta, quelle di bordo due volte, quelle interne quattro volte.

\begin{notebox}[HOG per il rilevamento di pedoni]
Per una patch $64\times128$ con celle $8\times8$ e blocchi $2\times2$ (step 8
pixel), si ottengono $7\times15=105$ blocchi, ciascuno con $4\times9=36$ valori.
Il descrittore HOG finale è un \textbf{vettore a 3780 dimensioni}:
\[
  7\times15\times36 = 3780
\]
Questo vettore può essere alimentato in un classificatore lineare (es.\ SVM).
La visualizzazione del descrittore HOG cattura chiaramente la forma della persona,
specialmente intorno al tronco e alle gambe.
\end{notebox}

\begin{propbox}[Proprietà di HOG]
HOG standard non ha assegnazione di orientazione dominante: è adatto a oggetti
con posa canonica (pedoni in piedi) e \textbf{non è invariante alla rotazione}.
\`E robusto ai cambiamenti locali di illuminazione grazie alla normalizzazione
per blocco, ma non invariante alla scala.
\end{propbox}

% =============================================================================
\chapter{Scale-Invariant Feature Transform (SIFT)}
% =============================================================================

Il \textit{Scale-Invariant Feature Transform} (\textbf{SIFT}) è un algoritmo di
riferimento per \textbf{rilevare} e \textbf{descrivere} feature locali nelle
immagini, progettato per essere robusto a variazioni di scala, rotazione e
illuminazione. SIFT è sia un rilevatore che un descrittore; ogni keypoint è
rappresentato come $(x,y,s,\theta,\mathbf{d})$.

\section{Rilevatore SIFT}

\subsection{Costruzione dello scale-space}

SIFT usa una piramide \textbf{Difference of Gaussians (DoG)}, approssimazione
efficiente del Laplaciano del Gaussiano (LoG):
\[
  \mathrm{DoG}(x,y,\sigma) = G(x,y,k\sigma)-G(x,y,\sigma)
\]
Sia DoG che LoG sono equivarianti alla rotazione.

Una piramide Gaussiana viene costruita in \textbf{ottave}. Ogni ottava contiene
immagini della stessa risoluzione, progressivamente sfocate. Per passare
all'ottava successiva si sottocampiona di un fattore 2. Le immagini DoG si
calcolano sottraendo livelli adiacenti all'interno di ogni ottava.

\subsection{Rilevamento degli estremi nello scale-space}

Un pixel è un keypoint potenziale se è un estremo locale rispetto ai suoi
\textbf{8 vicini spaziali} nella stessa immagine DoG e ai \textbf{9 vicini}
nelle immagini DoG sopra e sotto, per un totale di \textbf{26 vicini}.

\textbf{Rimozione dei keypoint instabili:}
\begin{itemize}
  \item \textbf{Basso contrasto}: risposta DoG debole, sensibile al rumore.
        Si scartano keypoint con risposta sotto una soglia.
  \item \textbf{Su bordi}: localizzazione ambigua lungo la direzione del bordo.
        La risposta DoG cambia fortemente in una sola direzione (trasversale al
        bordo) $\Rightarrow$ non distinti per il matching. Su un buon keypoint
        (blob o angolo) la risposta cambia significativamente in tutte le direzioni.
\end{itemize}

\begin{propbox}[Eliminazione dei keypoint sui bordi (Hessian Ratio Test)]
Per scartare i punti di interesse situati lungo i bordi rettilinei (dove la localizzazione è ambigua lungo la direzione del bordo), SIFT calcola la matrice Hessiana 2D $\mathbf{H}$ della risposta DoG nel punto candidato:
\[
  \mathbf{H} = \begin{bmatrix} D_{xx} & D_{xy} \\ D_{xy} & D_{yy} \end{bmatrix}
\]
Il rapporto tra i suoi autovalori $\alpha, \beta$ (curvature principali) viene verificato calcolando la traccia e il determinante senza calcolare gli autovalori espliciti:
\[
  \frac{\trace(\mathbf{H})^2}{\det(\mathbf{H})} = \frac{(\alpha + \beta)^2}{\alpha \beta} = \frac{(r+1)^2}{r}
\]
Se $\frac{\trace(\mathbf{H})^2}{\det(\mathbf{H})} \ge \frac{(r+1)^2}{r}$ (con soglia tipica $r=10$), significa che il punto si trova su un bordo e viene eliminato.
\end{propbox}

\subsection{Assegnazione dell'orientazione}

Per garantire l'invarianza alla rotazione si assegna un'\textbf{orientazione
canonica} a ogni keypoint:

\begin{enumerate}
  \item \textbf{Gradiente del vicinato}: si calcolano magnitudine e orientazione
        per i pixel nell'intorno del keypoint nell'immagine Gaussiana
        corrispondente.
  \item \textbf{Istogramma di orientazione}: istogramma a \textbf{36 bin}
        ($360^{\circ}$, bin da $10^{\circ}$); ogni pixel vota pesato dalla magnitudine e da
        una finestra Gaussiana centrata sul keypoint.
  \item \textbf{Orientazione/i dominante/i}: il picco più alto definisce
        l'orientazione primaria. Qualsiasi picco superiore all'\textbf{80\%}
        del principale genera un nuovo keypoint con stessa posizione e scala.
\end{enumerate}

\section{Descrittore SIFT}

Assegnati posizione $(x,y)$, scala $\sigma$ e orientazione $\theta$ al keypoint:

\begin{enumerate}
  \item Si prende una finestra $16\times16$ dall'immagine Gaussiana sfocata alla
        scala caratteristica.
  \item Si \textbf{ruotano} coordinate e gradienti secondo l'orientazione canonica
        $\Rightarrow$ \textbf{invarianza alla rotazione}.
  \item La finestra $16\times16$ viene divisa in una griglia $4\times4$ di celle
        ($4\times4$ pixel ciascuna, 16 celle in totale). Per ogni cella si calcola
        un \textbf{istogramma di orientazione a 8 bin}.
  \item Si \textbf{concatenano} i 16 istogrammi in un vettore a
        $\mathbf{128}$ \textbf{dimensioni} ($16\times8=128$).
  \item Si \textbf{normalizza} a lunghezza unitaria $\Rightarrow$ robustezza ai
        cambiamenti di illuminazione.
\end{enumerate}

\begin{propbox}[Proprietà di SIFT]
Il descrittore SIFT è invariante a:
\begin{itemize}
  \item \textbf{Scala}: rilevamento in piramide di scale-space.
  \item \textbf{Rotazione}: normalizzazione dell'orientazione.
  \item \textbf{Illuminazione}: normalizzazione del vettore descrittore.
\end{itemize}
\`E anche robusto alle trasformazioni affini e al rumore, ma non completamente
invariante a grandi cambiamenti di punto di vista.
\end{propbox}

\section{Matching di feature}

\subsection{Ratio test di Lowe}

Per ogni feature $f_1$ in $I_1$ si trovano i due vicini più prossimi in $I_2$:
$f_2$ (best match) e $f'_2$ (second-best). Si calcola il ratio:
\[
  \mathrm{ratio} = \frac{\|f_1-f_2\|}{\|f_1-f'_2\|}
\]
Il match è accettato se ratio $<$ soglia (tipicamente 0.75--0.8). Un ratio piccolo
indica un match non ambiguo: il miglior match è significativamente migliore del
secondo migliore.

\section{Valutazione del matching}

Variando la soglia si ottiene la \textbf{ROC curve}:
\[
  \mathrm{TPR}=\frac{\#\mathrm{TP}}{\#\text{match corretti totali}}, \qquad
  \mathrm{FPR}=\frac{\#\mathrm{FP}}{\#\text{match scorretti totali}}
\]

\begin{notebox}[Metriche informative]
La semplice accuracy è scadente per il feature matching: i veri negativi
(match non corrispondenti) sono astronomicamente numerosi e dominano il calcolo.
Preferire:
\[
  \mathrm{Precision}=\frac{TP}{TP+FP},\quad
  \mathrm{Recall}=\frac{TP}{TP+FN},\quad
  F_1=2\frac{P\cdot R}{P+R}
\]
L'\textbf{AUC} (Area Under the ROC Curve) riassume la performance complessiva:
1.0 = classificatore perfetto, 0.5 = random.
\end{notebox}

\section{SURF: Speeded-Up Robust Features}

SURF è un'alternativa più veloce a SIFT con qualità comparabile.

\textbf{Rilevamento (Hessiano):} SURF rileva blob dove $\det(\mathbf{H})$ è
massimo, con:
\[
  \mathbf{H}(\mathbf{x},\sigma)=
  \begin{bmatrix}L_{xx}&L_{xy}\\L_{xy}&L_{yy}\end{bmatrix}, \qquad
  \det(H_{\mathrm{approx}})=D_{xx}D_{yy}-(wD_{xy})^2, \quad w=0.9
\]
Le derivate Gaussiane del secondo ordine vengono approssimate con \textbf{filtri
box} calcolati tramite \textbf{immagini integrali} per massima efficienza.

\textbf{Descrittore:} finestra $20s\times20s$ allineata all'orientazione dominante,
divisa in griglia $4\times4$. Per ogni cella: $(\sum dx, \sum dy, \sum|dx|,
\sum|dy|)$ $\Rightarrow$ \textbf{vettore a 64 dimensioni}.

\section{Descrittori binari di feature}

Alternative meno costose a SIFT che generano \textbf{stringhe binarie brevi},
molto veloci da calcolare e confrontare.

Per ogni coppia di pixel $(s_1,s_2)$ nella patch:
\[
  b=\begin{cases}1&\text{se }I(s_1)<I(s_2)\\0&\text{altrimenti}\end{cases}
\]

\begin{propbox}[Vantaggi]
Compatti (es.\ 256 bit per ORB); veloci da calcolare (semplici confronti di
intensità); veloci da confrontare tramite \textbf{distanza di Hamming}:
$d_H(B_1,B_2)=\mathrm{sum}(\mathrm{xor}(B_1,B_2))$.
\end{propbox}

\textbf{ORB (Oriented FAST and Rotated BRIEF):} migliora BRIEF aggiungendo
invarianza alla rotazione. Stima il centro di massa (CoM) e ruota di $\theta$ le
256 coppie apprese da training (coppie non correlate ad alta varianza) prima di
eseguire i confronti.

\section{Deep local features}

\textbf{SuperPoint} apprende detection e description tramite dati sintetici e
self-supervised learning. \textbf{SuperGlue} apprende il matching con graph
neural networks.

% =============================================================================
\chapter{Trasformazioni e allineamento}
% =============================================================================

\section{Image warping}

\begin{itemize}
  \item \textbf{Image filtering}: cambia i valori dei pixel (range)
        $\Rightarrow g(x)=h(f(x))$.
  \item \textbf{Image warping}: cambia le coordinate dei pixel (dominio)
        $\Rightarrow g(x)=f(h(x))$, dove $h$ è una trasformazione di coordinate.
\end{itemize}

Ci si concentra sul \textbf{global warping}: una singola trasformazione
$\mathbf{p}'=T(\mathbf{p})$ applicata a tutti i punti $\mathbf{p}=(x,y)$.

\section{Tipi di trasformazioni 2D}

\begin{center}
\begin{tabular}{lccl}
\toprule
\textbf{Trasformazione}&\textbf{DOF}&\textbf{Min. corrisp.}&\textbf{Preserva}\\
\midrule
Traslazione    & 2 & 1 & orientazione, lunghezze, angoli, parallelismo\\
Euclidea/rigid & 3 & --- & lunghezze, angoli, parallelismo\\
Similarità     & 4 & --- & angoli, parallelismo, rapporti\\
Affine         & 6 & 3 & linee, parallelismo, rapporti su parallele\\
Proiettiva     & 8 & 4 & linee rette\\
\bottomrule
\end{tabular}
\end{center}

\subsection{Trasformazioni lineari}

\[
  \mathbf{p}'=\mathbf{T}\mathbf{p}\implies
  \begin{bmatrix}x'\\y'\end{bmatrix}=
  \begin{bmatrix}a&b\\c&d\end{bmatrix}\begin{bmatrix}x\\y\end{bmatrix}
\]

Esempi:
\begin{itemize}
  \item Scala: $\mathbf{S}=\begin{bmatrix}s_x&0\\0&s_y\end{bmatrix}$
  \item Rotazione: $\mathbf{R}=\begin{bmatrix}\cos\theta&-\sin\theta\\\sin\theta&\cos\theta\end{bmatrix}$
  \item Shear: $\mathbf{Sh}=\begin{bmatrix}1&sh_x\\sh_y&1\end{bmatrix}$
\end{itemize}

\begin{propbox}[Proprietà]
L'origine mappa nell'origine; le linee mappano in linee; il parallelismo è
preservato; i rapporti di lunghezze lungo una linea sono preservati; chiuse sotto
composizione.
\end{propbox}

\begin{warnbox}[Limitazione delle matrici $2\times2$]
La traslazione non è lineare su coordinate 2D e non può essere rappresentata da
una matrice $2\times2$.
\end{warnbox}

\subsection{Trasformazioni affini}

Con \textbf{coordinate omogenee} $(x,y)\to(x,y,1)$ si includono le traslazioni:
\[
  \begin{bmatrix}x'\\y'\\1\end{bmatrix}=
  \begin{bmatrix}a&b&c\\d&e&f\\0&0&1\end{bmatrix}
  \begin{bmatrix}x\\y\\1\end{bmatrix}
\]

\begin{notebox}[Da coordinate omogenee a 2D]
$(x,y,w)\to(x/w,\,y/w)$ (normalizzazione). I punti con $w=0$ sono punti
all'infinito (utili per i punti di fuga).
\end{notebox}

La trasformazione affine mappa linee in linee, preserva il parallelismo e i
rapporti lungo linee parallele. L'origine non mappa necessariamente nell'origine.

\subsection{Trasformazioni proiettive (omografie)}

Matrice $3\times3$ $\mathbf{H}$ definita a meno di un fattore di scala (8 DOF):
\[
  \begin{bmatrix}x'\\y'\\w'\end{bmatrix}=
  \underbrace{\begin{bmatrix}a&b&c\\d&e&f\\g&h&1\end{bmatrix}}_{\mathbf{H}}
  \begin{bmatrix}x\\y\\w\end{bmatrix},
  \qquad \text{coordinate Cartesiane: }(x'/w',\,y'/w')
\]

\begin{propbox}[Proprietà]
Le linee rette mappano in linee rette. Le linee parallele non rimangono
necessariamente parallele (convergono a punti di fuga). I rapporti non sono
preservati.
\end{propbox}

\begin{notebox}[Applicazione]
Un'omografia mappa punti tra due immagini diverse dello stesso piano 3D.
\`E la trasformazione che relaziona due viste della stessa superficie planare
nel modello a camera pinhole.
\end{notebox}

\section{Calcolo delle trasformazioni}

\subsection{Sistema lineare e minimi quadrati}

Con più della minima quantità di corrispondenze si ottiene un sistema
\textbf{sovradeterminato} $\mathbf{At}=\mathbf{b}$, risolto con i
\textbf{minimi quadrati lineari}:
\[
  \mathbf{t}=(\mathbf{A}^T\mathbf{A})^{-1}\mathbf{A}^T\mathbf{b}=\mathbf{A}^\#\mathbf{b}
\]

\subsection{Calcolo di un'omografia}

Per una corrispondenza $(x_i,y_i)\to(x'_i,y'_i)$, riorganizzando si ottiene
$\mathbf{Ah}=\mathbf{0}$ con due righe per corrispondenza:
\[
  \begin{bmatrix}
    x_i&y_i&1&0&0&0&-x'_ix_i&-x'_iy_i&-x'_i\\
    0&0&0&x_i&y_i&1&-y'_ix_i&-y'_iy_i&-y'_i
  \end{bmatrix}\mathbf{h}=\mathbf{0}
\]

Con 4 corrispondenze (8 equazioni) si risolve il sistema omogeneo:
minimizzare $\|\mathbf{Ah}\|^2$ soggetto a $\|\mathbf{h}\|=1$.

\begin{notebox}[Soluzione SVD]
$\mathbf{h}$ è l'autovettore di $\mathbf{A}^T\mathbf{A}$ per il minimo
autovalore. Si usa la SVD di $\mathbf{A}$: $\mathbf{h}$ è l'ultima colonna di
$\mathbf{V}$.
\end{notebox}

\section{Stima robusta con RANSAC}

I minimi quadrati sono sensibili agli \textbf{outlier} (match scorretti).
\textbf{RANSAC} (RANdom SAmple Consensus) identifica il più grande insieme
di inlier:

\begin{enumerate}
  \item \textbf{Campiona} $s$ punti (2 per lineare, 3 per affine, 4 per omografia).
  \item \textbf{Calcola} il modello dal campione minimale.
  \item \textbf{Conta} gli inlier entro tolleranza $\varepsilon$.
  \item \textbf{Ripeti} per $N$ iterazioni; conserva il modello con più inlier.
  \item \textbf{Ricalcola} il modello con tutti gli inlier tramite minimi quadrati.
\end{enumerate}

\begin{defbox}[Numero teorico di iterazioni in RANSAC]
Il numero minimo di iterazioni $N$ necessario a RANSAC per garantire con probabilità $p$ (es.\ 99\%) di pescare almeno un sottoinsieme di $s$ punti composto esclusivamente da inlier (con un tasso stimato di outlier $e$) è:
\[
  N \ge \frac{\log(1 - p)}{\log\bigl(1 - (1 - e)^s\bigr)}
\]
\textit{Implicazione d'esame}: All'aumentare della percentuale di outlier ($e \to 1$), il denominatore tende a zero e il numero di iterazioni $N$ necessarie cresce in modo \textbf{esponenziale}.
\end{defbox}

\section{Implementazione del warping e blending}

\subsection{Forward vs.\ inverse warping}

\begin{itemize}
  \item \textbf{Forward}: per ogni pixel $(x,y)$ sorgente si calcola
        $(x',y')=T(x,y)$ e si copia il valore. \textit{Problema}: buchi e
        sovrapposizioni nell'output.
  \item \textbf{Inverse}: per ogni pixel $(x',y')$ destinazione si calcola
        $(x,y)=T^{-1}(x',y')$ e si campiona il colore. \textbf{Garantisce ogni
        pixel riempito}. Poiché $(x,y)$ non è in genere intero, serve
        \textbf{interpolazione} (nearest neighbor, bilineare, bicubica).
\end{itemize}

\subsection{Blending}

\begin{itemize}
  \item \textbf{Feathering}: media pesata (alpha blending) con transizione
        graduale nelle zone sovrapposte. Può produrre ghosting se l'allineamento
        è imperfetto.
  \item \textbf{Pyramid blending}: si veda Capitolo~3.
  \item \textbf{Poisson blending}: metodo nel dominio del gradiente che risolve
        un'equazione di Poisson per preservare i gradienti nella regione incollata.
\end{itemize}

\section{Punti di fuga}

Un \textbf{punto di fuga} è la proiezione di una direzione 3D. Le linee 3D
parallele con la stessa direzione convergono in esso. Per due rette $l_1,l_2$
in coordinate omogenee: $v=l_1\times l_2$.

% =============================================================================
\chapter{Riconoscimento}
% =============================================================================

\section{Compiti di riconoscimento}

\begin{center}
\begin{tabular}{ll}
\toprule
\textbf{Compito}&\textbf{Output}\\
\midrule
Image classification  & Un'etichetta di classe per l'intera immagine\\
Semantic segmentation & Un'etichetta semantica per ogni pixel\\
Object detection      & Bounding box ed etichette di classe\\
Instance segmentation & Etichetta semantica e identità di istanza per i pixel degli oggetti\\
\bottomrule
\end{tabular}
\end{center}

\section{Classificazione di immagini}

Si usano approcci di ML con dataset etichettati (es.\ ImageNet). Il
\textbf{transfer learning} consente il pre-training su un dataset grande e il
fine-tuning su uno più piccolo.

\textbf{Baseline nearest neighbor}: si assegna la classe dell'immagine di training
più vicina. Le distanze pixel-wise sono però poco informative a causa di shift,
illuminazione o deformazioni.

\subsection{Bag-of-Words (BoW)}

BoW raggiunge l'\textbf{invarianza spaziale} rappresentando immagini come
istogrammi di feature locali, ignorando la loro disposizione spaziale.

\textbf{Pipeline:}
\begin{enumerate}
  \item Rilevare e descrivere feature (es.\ SIFT).
  \item Raggruppare feature con k-means per formare il vocabolario visivo.
  \item Assegnare ogni descrittore alla parola visiva più vicina.
  \item Rappresentare ogni immagine con l'istogramma di frequenza delle parole.
  \item Addestrare un classificatore (es.\ SVM) sugli istogrammi.
\end{enumerate}

\section{Reti Neurali Convoluzionali (CNN)}

Tre tipi principali di strati:
\begin{itemize}
  \item \textbf{Convoluzionali}: filtri apprendibili con \textbf{weight sharing}.
        Equivarianti alla traslazione.
  \item \textbf{Pooling}: riducono la risoluzione spaziale e aumentano il
        \textbf{campo ricettivo}. Solitamente senza parametri apprendibili.
  \item \textbf{FC}: classificazione finale.
\end{itemize}

\subsection{Output e funzione di loss}

Funzione \textbf{softmax}:
\[
  p_c=\frac{e^{s_c}}{\sum_j e^{s_j}}
\]
Loss \textbf{cross-entropy}:
\[
  L=-\log p_y
\]

\subsection{Architetture principali}

\begin{center}
\begin{tabular}{ll}
\toprule
\textbf{Architettura}&\textbf{Idea chiave}\\
\midrule
LeNet-5            & Prima CNN per cifre\\
AlexNet            & ImageNet breakthrough, ReLU, dropout, data augmentation\\
VGG                & Stack profondi di conv.\ $3\times3$\\
GoogLeNet/Inception& Moduli convoluzionali multi-scala, meno parametri\\
ResNet             & Connessioni residuali per reti molto profonde\\
\bottomrule
\end{tabular}
\end{center}

\textbf{ResNet} mitiga il \textbf{vanishing gradient} con connessioni residuali
che permettono al gradiente di bypassare gli strati.

\section{Segmentazione semantica}

Usa l'\textbf{architettura encoder-decoder}:
\begin{itemize}
  \item \textit{Encoder}: downsampling progressivo per feature di alto livello.
  \item \textit{Decoder}: upsampling progressivo per produrre la mappa di
        segmentazione.
  \item \textbf{Skip connections}: combinano informazioni grossolane del decoder
        con quelle fini dell'encoder.
\end{itemize}

\textbf{Metriche}: IoU (Jaccard Index) $=\frac{|P\cap G|}{|P\cup G|}$,
accuratezza per pixel.

\textbf{Fully Convolutional Networks (FCN):} sostituiscono i layer FC con
operazioni convoluzionali producendo heatmap spaziali. Skip connections recuperano
il dettaglio spaziale.

\textbf{SegNet:} encoder-decoder con pooling indices per guidare l'unpooling.

\textbf{Convoluzione dilatata:} aumenta il campo ricettivo senza aumentare i
parametri e senza ridurre la risoluzione spaziale. Utile per predizione densa.

\textbf{U-Net:} encoder-decoder con skip connections tra livelli corrispondenti.
Le skip connections concatenano feature a bassa e alta risoluzione, migliorando
la localizzazione. Standard de facto per segmentazione biomedicale.

\section{Rilevamento di oggetti}

\textbf{mAP} (mean Average Precision): metrica principale calcolata dalle curve
precisione-recall, con soglia IoU $\geq0.5$ per i veri positivi.

\subsection{Rilevatori a due stadi}

\begin{itemize}
  \item \textbf{R-CNN (2014)}: Selective Search per $\sim$2k proposte; CNN su
        ciascuna indipendentemente. Lento.
  \item \textbf{Fast R-CNN (2015)}: CNN sull'intera immagine; \textbf{RoI Pooling}
        per feature a dimensione fissa da ogni proposta.
  \item \textbf{Faster R-CNN (2015)}: \textbf{Region Proposal Network (RPN)}
        apprendibile; predice objectness e raffinamenti rispetto ad
        \textbf{anchor box} predefiniti. Molto più veloce.
\end{itemize}

\subsection{Rilevatori a stadio singolo}

\textbf{YOLO} e \textbf{SSD} predicono bounding box e probabilità di classe
direttamente in un singolo passaggio. Generalmente più veloci.

\subsection{Feature Pyramid Networks (FPN)}

Costruiscono una \textbf{piramide di feature multi-scala} con percorso top-down
e connessioni laterali, permettendo il rilevamento a multiple scale.

\subsection{Sliding window e HOG}

Approccio classico: finestra scorrevole + HOG + classificatore SVM.
Non-maximum suppression rimuove i rilevamenti duplicati. Semplice ma lento.

\section{Estensioni}

\textbf{Mask R-CNN:} estende Faster R-CNN con un ramo di segmentazione per ogni
oggetto rilevato. Sostituisce RoI Pooling con \textbf{RoIAlign} per preservare
l'allineamento spaziale.

\begin{notebox}[RoIPooling vs.\ RoIAlign: Il problema della Quantizzazione]
In Faster R-CNN, il \textbf{RoIPooling} esegue un arrotondamento (quantizzazione) sulle coordinate decimali della regione d'interesse per adattarle alla griglia discreta dei pixel della feature map. Questo errore di disallineamento spaziale ($\approx 0.5$ pixel) è trascurabile per le Bounding Box, ma distruttivo per le maschere di segmentazione al pixel.

\textbf{Mask R-CNN} sostituisce RoIPooling con \textbf{RoIAlign}: evita qualsiasi quantizzazione e calcola il valore esatto delle caratteristiche in punti a coordinate decimali utilizzando l'\textbf{interpolazione bilineare} a partire da 4 punti di campionamento, preservando un perfetto allineamento spaziale.
\end{notebox}

\textbf{Segmentazione panoptica:} unifica segmentazione semantica (``stuff'':
cielo, strada) e di istanze (``things'': auto, persone).

% =============================================================================
\chapter{Flusso ottico}
% =============================================================================

\section{Introduzione}

\begin{defbox}[Flusso ottico]
Il \textbf{flusso ottico} è il moto apparente di oggetti, superfici e bordi in
una scena visiva causato dal moto relativo tra osservatore e scena.
\`E un \textbf{campo vettoriale 2D} dove ogni vettore è uno spostamento di pixel
tra fotogrammi successivi.
\end{defbox}

Visualizzazione tipica: codifica HSV, dove la tonalità indica la direzione e la
luminosità la magnitudine.

\begin{notebox}[Campo di moto vs.\ flusso ottico]
\begin{itemize}
  \item \textbf{Campo di moto}: vera proiezione 2D del moto 3D sul piano immagine.
  \item \textbf{Flusso ottico}: moto apparente dei pattern di luminosità.
        Possono differire: sfera Lambertiana ruotante sotto illuminazione fissa
        $\to$ campo di moto ma flusso ottico nullo; sfera speculare ferma con
        sorgente di luce mobile $\to$ flusso ottico ma campo di moto nullo.
\end{itemize}
\end{notebox}

\begin{notebox}[Stereo vs.\ flusso ottico]
\begin{center}
\begin{tabular}{ll}
\toprule
\textbf{Stereo}&\textbf{Flusso ottico}\\
\midrule
Due immagini allo stesso tempo & Due immagini in tempi diversi\\
Ricerca 1D lungo epipolare & Stima moto 2D (tutte le direzioni)\\
Usa geometria epipolare & Nessun vincolo epipolare diretto\\
\bottomrule
\end{tabular}
\end{center}
\end{notebox}

\section{Il problema dell'apertura}

Osservando un bordo in movimento in una piccola patch locale, solo il moto
perpendicolare al bordo può essere determinato (\textbf{flusso normale}); la
componente parallela rimane ambigua. Per evitare il problema si necessita di
patch con gradienti in direzioni diverse (angoli, texture).

\section{Stima del flusso ottico}

\textbf{Assunzioni:}
\begin{itemize}
  \item \textbf{Costanza della luminosità}: $I(x(t),y(t),t)$ costante nel tempo.
  \item \textbf{Piccolo moto}: i punti non si spostano troppo tra fotogrammi.
  \item \textbf{Coerenza spaziale}: i punti vicini si muovono in modo simile.
\end{itemize}

\subsection{Equazione di costanza della luminosità}

Dalla costanza della luminosità e dall'espansione di Taylor al primo ordine:
\[
  I_x u + I_y v + I_t = 0
\]
dove $I_x,I_y$ sono i gradienti spaziali e $I_t$ è il gradiente temporale.

\begin{warnbox}[Problema sottodeterminato]
Un'equazione per pixel, due incognite $(u,v)$: le soluzioni giacciono su una
retta nello spazio $(u,v)$ $\Rightarrow$ problema sottodeterminato per singoli
pixel (apertura problem).
\end{warnbox}

\section{Lucas-Kanade}

Assume il flusso \textbf{costante} in una piccola patch $N\times N$.
Sistema sovradeterminato $\mathbf{Ax}=\mathbf{b}$:
\[
  \underbrace{\begin{bmatrix}I_x(\mathbf{p}_1)&I_y(\mathbf{p}_1)\\\vdots&\vdots\\I_x(\mathbf{p}_{N^2})&I_y(\mathbf{p}_{N^2})\end{bmatrix}}_{\mathbf{A}}
  \begin{bmatrix}u\\v\end{bmatrix}=
  \underbrace{\begin{bmatrix}-I_t(\mathbf{p}_1)\\\vdots\\-I_t(\mathbf{p}_{N^2})\end{bmatrix}}_{\mathbf{b}}
\]

Soluzione minimi quadrati: $\mathbf{x}=(\mathbf{A}^T\mathbf{A})^{-1}\mathbf{A}^T\mathbf{b}$.

La matrice $\mathbf{A}^T\mathbf{A}$ è la \textbf{matrice dei secondi momenti}
(come in Harris):
\[
  \mathbf{A}^T\mathbf{A}=\sum_{\mathbf{p}}\begin{bmatrix}I_x^2&I_xI_y\\I_xI_y&I_y^2\end{bmatrix}
\]

\begin{propbox}[Quando funziona Lucas-Kanade?]
\begin{center}
\begin{tabular}{lll}
\toprule
\textbf{Regione}&\textbf{Autovalori $\lambda_1,\lambda_2$}&\textbf{Flusso}\\
\midrule
Piatta    & Entrambi piccoli & Non affidabile\\
Bordo     & Uno grande, uno piccolo & Ambiguo\\
Angolo/texture & Entrambi grandi & Affidabile\\
\bottomrule
\end{tabular}
\end{center}
\end{propbox}

\subsection{Gestione dei grandi moti: coarse-to-fine}

L'assunzione di piccolo moto fallisce per spostamenti grandi. Si usa la strategia
\textbf{coarse-to-fine} con piramidi Gaussiane:

\begin{enumerate}
  \item Costruire piramidi Gaussiane per entrambi i fotogrammi.
  \item Stimare $(u,v)^{(0)}$ al livello più grossolano (grandi moti appaiono
        piccoli).
  \item Iterativamente: upsample del flusso $\to$ warp del fotogramma $t+1$
        verso $t$ (inverse warping) $\to$ calcolo del residuo $\Delta(u,v)$
        $\to$ accumulo.
  \item Ripetere fino alla risoluzione piena.
\end{enumerate}

\section{Horn-Schunck}

Assume il flusso \textbf{regolare} sull'intera immagine. Minimizza:
\[
  E(u,v)=
  \underbrace{\iint(I_xu+I_yv+I_t)^2\,dxdy}_{\text{termine dati}}+
  \lambda\underbrace{\iint(\|\nabla u\|^2+\|\nabla v\|^2)\,dxdy}_{\text{regolarizzazione}}
\]

\begin{warnbox}[Oversmoothing]
La penalità quadratica penalizza troppo le discontinuità di flusso ai confini
degli oggetti, causando \textbf{oversmoothing}. Si risolve con funzioni di loss
\textbf{robuste} (es.\ penalità Student-t) meno sensibili ai grandi errori.
\end{warnbox}

\section{Deep learning per il flusso ottico}

\textbf{FlowNet:} prima rete end-to-end supervisionata per flusso ottico.
Architettura encoder-decoder su due fotogrammi consecutivi.
\begin{itemize}
  \item FlowNetSimple: stack dei due frame in input a una singola rete.
  \item FlowNetCorr: due stream con correlation layer per il matching.
\end{itemize}

\textbf{End-Point Error (EPE):}
$\mathrm{EPE}=\sqrt{(u-\hat{u})^2+(v-\hat{v})^2}$

\textbf{FlowNet2:} impila più moduli FlowNet con warping tra i moduli.

\textbf{PWC-Net} (\textbf{P}yramid, \textbf{W}arping, \textbf{C}ost volume):
piramide di feature, stima coarse-to-fine, warping delle feature del secondo frame
verso il primo, cost volume per il matching. Combina principi classici e DL.

\textbf{Flusso non supervisionato:} apprendimento senza ground-truth tramite
loss di ricostruzione fotometrica e di smoothness.

% =============================================================================
\chapter{Calibrazione della camera}
% =============================================================================

\section{Il modello della camera}

Una camera esegue una \textbf{trasformazione proiettiva} che mappa il mondo 3D
su un piano immagine 2D. La \textbf{calibrazione} è il processo di stima dei
parametri di questa trasformazione.

\subsection{Pinhole e lente sottile}

Il modello più semplice è la \textbf{camera pinhole}: una piccola apertura lascia
passare un solo raggio per punto, producendo un'immagine invertita senza sfocatura.

\begin{notebox}[Trade-off del pinhole]
Apertura piccola $\to$ meno blur ma effetti di diffrazione e poca luce
(lunga esposizione). Formula approssimata per il diametro ideale:
$d\approx2\sqrt{f\lambda}$.
\end{notebox}

\textbf{Legge della lente sottile:}
\[
  \frac{1}{i}+\frac{1}{o}=\frac{1}{f}
\]
dove $o$ è la distanza oggetto, $i$ la distanza immagine, $f$ la lunghezza focale.
Solo gli oggetti a una profondità specifica sono a fuoco.

\subsection{Proiezione prospettica}

Relazione tra punto 3D $\mathbf{x}_c=(x_c,y_c,z_c)$ nel frame camera e
proiezione 2D $(x_i,y_i)$:
\[
  x_i=f\frac{x_c}{z_c}, \qquad y_i=f\frac{y_c}{z_c}
\]

\subsection{Dal piano immagine ai pixel}

\[
  u=f_x\frac{x_c}{z_c}+o_x, \qquad v=f_y\frac{y_c}{z_c}+o_y
\]
dove $f_x=m_xf$, $f_y=m_yf$ (lunghezze focali in pixel) e $(o_x,o_y)$ è il
punto principale.

\subsection{La matrice della camera}

\[
  \tilde{\mathbf{u}}=\mathbf{P}\tilde{\mathbf{x}}_w=\mathbf{K}[\mathbf{R}|\mathbf{t}]\tilde{\mathbf{x}}_w
\]

\textbf{Matrice intrinseca} $\mathbf{K}$ (parametri intrinseci $f_x,f_y,o_x,o_y$):
\[
  \mathbf{K}=\begin{bmatrix}f_x&0&o_x\\0&f_y&o_y\\0&0&1\end{bmatrix}
\]

\textbf{Matrice estrinseca} $[\mathbf{R}|\mathbf{t}]$ (parametri estrinseci):
\[
  [\mathbf{R}|\mathbf{t}]=\begin{bmatrix}r_{11}&r_{12}&r_{13}&t_x\\r_{21}&r_{22}&r_{23}&t_y\\r_{31}&r_{32}&r_{33}&t_z\end{bmatrix}
\]

\begin{notebox}[Distorsione dell'obiettivo]
Gli obiettivi reali introducono distorsioni radiali e tangenziali:
\[
  \begin{pmatrix}x'\\y'\end{pmatrix}=
  (1+\kappa_1r^2+\kappa_2r^4)\begin{pmatrix}x\\y\end{pmatrix}+
  \begin{pmatrix}2\kappa_3xy+\kappa_4(r^2+2x^2)\\2\kappa_4xy+\kappa_3(r^2+2y^2)\end{pmatrix}
\]
con $r^2=x^2+y^2$. Poi: $(u,v)=(f_x x'+c_x,\,f_y y'+c_y)$.
\end{notebox}

\section{Calibrazione della camera}

Si stima $\mathbf{P}$ osservando un oggetto di geometria nota (es.\ scacchiera)
con almeno 6 coppie di punti 3D noti e corrispondenti punti immagine (11 incognite
perché $\mathbf{P}$ è definita a meno di scala).

\begin{defbox}[Espansione delle equazioni DLT per singolo punto]
Ogni corrispondenza tra punto 3D del mondo $\mathbf{X}_i = (x_i, y_i, z_i, 1)^T$ e pixel 2D $\mathbf{u}_i = (u_i, v_i, 1)^T$ soddisfa $\mathbf{u}_i \sim \mathbf{P} \mathbf{X}_i$. Eliminando lo scalare omogeneo tramite divisione per la terza riga, si ricavano due equazioni lineari indipendenti per ciascun punto:
\[
  \begin{aligned}
    u_i (p_{31} x_i + p_{32} y_i + p_{33} z_i + p_{34}) - (p_{11} x_i + p_{12} y_i + p_{13} z_i + p_{14}) &= 0 \\
    v_i (p_{31} x_i + p_{32} y_i + p_{33} z_i + p_{34}) - (p_{21} x_i + p_{22} y_i + p_{23} z_i + p_{24}) &= 0
  \end{aligned}
\]
Impilando queste due righe per $N \ge 6$ punti, si compone la matrice dei dati $\mathbf{A}_{2N \times 12}$ per il sistema lineare omogeneo $\mathbf{A}\mathbf{p} = \mathbf{0}$.
\end{defbox}

Ogni corrispondenza dà 2 equazioni. Sistema omogeneo $\mathbf{Ap}=\mathbf{0}$:
minimizzare $\|\mathbf{Ap}\|^2$ soggetto a $\|\mathbf{p}\|=1$.
Soluzione: autovettore di $\mathbf{A}^T\mathbf{A}$ per il minimo autovalore.

\subsection{Estrazione dei parametri}

Tramite \textbf{fattorizzazione RQ}:
\[
  \mathbf{P}=[\mathbf{M}|\mathbf{p}_4]\implies
  \mathbf{M}=\mathbf{KR},\;\mathbf{p}_4=\mathbf{Kt}\implies
  \mathbf{K}=\mathrm{RQ}(\mathbf{M}),\;
  \mathbf{R}=\mathbf{K}^{-1}\mathbf{M},\;
  \mathbf{t}=\mathbf{K}^{-1}\mathbf{p}_4
\]

\section{Punti di fuga}

Un \textbf{punto di fuga} è la proiezione di una direzione 3D. Le linee 3D
parallele convergono in esso. Per due rette $l_1,l_2$: $v=l_1\times l_2$.

\section{Visione stereo}

Usa \textbf{due o più camere} per inferire profondità tramite triangolazione.

\subsection{Profondità dalla disparità}

Per stereo rettificato (camere parallele, baseline $b$):
\[
  z=\frac{b\,f_x}{d}=\frac{b\,f_x}{u_l-u_r}
\]
Alta disparità $d$ $\to$ punto vicino; bassa $\to$ lontano.

\subsection{Matching stereo}

\textbf{Matching a finestra:}
\begin{enumerate}
  \item Patch intorno al pixel in $I_l$.
  \item Confronto con patch lungo la stessa scanline in $I_r$.
  \item Best match con SSD/SAD/NCC; $d=u_l-u_r$.
\end{enumerate}

\begin{notebox}[Trade-off della finestra]
Finestra piccola: cattura dettagli fini, mappa rumorosa.
Finestra grande: robusta al rumore, sfoca le discontinuità.
Casi problematici: regioni senza texture, pattern ripetuti, superfici
speculari, discontinuità di profondità.
\end{notebox}

\subsection{Stereo come minimizzazione di energia}

\[
  E(d)=\underbrace{\sum_{(x,y)}E_d(x,y,d(x,y))}_{\text{dati}}+
       \lambda\underbrace{\sum_{(p,q)\in\mathcal{E}}E_s(d_p,d_q)}_{\text{smoothness}}
\]

Termini di smoothness: Potts ($\mathbf{1}[d_p\neq d_q]$) o L1 ($|d_p-d_q|$).
Ottimizzazione: programmazione dinamica per scanline (ma streaking) o graph cuts
(globale).

\section{Stereo attivo e luce strutturata}

Un proiettore proietta pattern di luce strutturata sulla scena, creando texture
artificiale e semplificando le corrispondenze. Base di sensori come Kinect e
Apple TrueDepth.

% =============================================================================
\chapter{Geometria epipolare}
% =============================================================================

\section{Introduzione}

\begin{defbox}[Geometria epipolare]
La \textbf{geometria epipolare} mette in relazione le proiezioni 2D
$(\mathbf{x},\mathbf{x}')$ di un punto 3D $\mathbf{X}$ da due viste diverse,
fornendo un vincolo geometrico per trovare corrispondenze e pose relative.
\end{defbox}

La geometria multi-vista relaziona tre componenti: moto della camera,
corrispondenze tra immagini e struttura 3D.
\begin{itemize}
  \item \textbf{Calibrazione / pose estimation}: da corrispondenze 3D-2D note.
  \item \textbf{Triangolazione}: da camera note e corrispondenze 2D-2D.
  \item \textbf{Structure from Motion / SLAM}: stima sia moto che struttura.
\end{itemize}

\textbf{Back-projection:} da una singola immagine, un pixel non determina un
punto 3D unico ma un raggio retroproiettato. Serve una seconda vista.

\textbf{Definizioni:}
\begin{itemize}
  \item \textbf{Baseline}: segmento tra i centri $\mathbf{O},\mathbf{O}'$.
  \item \textbf{Piano epipolare}: piano definito da $\mathbf{X},\mathbf{O},\mathbf{O}'$.
  \item \textbf{Epipoli} $(\mathbf{e},\mathbf{e}')$: intersezione della baseline
        con i piani immagine. $\mathbf{e}$ = proiezione di $\mathbf{O}'$,
        $\mathbf{e}'$ = proiezione di $\mathbf{O}$.
  \item \textbf{Linee epipolari} $(\mathbf{l},\mathbf{l}')$: intersezione del
        piano epipolare con i piani immagine.
\end{itemize}

\textbf{Vincolo epipolare:} $\mathbf{x}'$ deve giacere su $\mathbf{l}'$
$\Rightarrow$ riduce la ricerca di corrispondenze da 2D a 1D.

\section{La matrice essenziale}

\begin{defbox}[Matrice essenziale]
$\mathbf{E}\in\R^{3\times3}$ codifica il vincolo epipolare per camere
\textbf{calibrate} ($\mathbf{K}$ nota):
\[
  \boxed{\mathbf{x}'^T\mathbf{E}\mathbf{x}=0}
\]
Linee: $\mathbf{l}'=\mathbf{E}\mathbf{x}$, $\mathbf{l}=\mathbf{E}^T\mathbf{x}'$.
Epipoli: $\mathbf{Ee}=\mathbf{0}$, $\mathbf{e}'^T\mathbf{E}=\mathbf{0}$.
\end{defbox}

\subsection{Derivazione}

Con $\mathbf{x}=\mathbf{K}^{-1}\mathbf{p}$, $\mathbf{x}'=\mathbf{R}(\mathbf{x}-\mathbf{t})$.
La coplanarità di $\mathbf{x}$, $\mathbf{t}$, $\mathbf{x}-\mathbf{t}$ implica:
\[
  (\mathbf{x}-\mathbf{t})^T(\mathbf{t}\times\mathbf{x})=0
  \implies \mathbf{x}'^T\mathbf{R}[\mathbf{t}]_\times\mathbf{x}=0
  \implies \mathbf{E}=\mathbf{R}[\mathbf{t}]_\times
\]

\begin{propbox}[Proprietà di $\mathbf{E}$]
\begin{itemize}
  \item Opera su coordinate normalizzate.
  \item Rango 2, $\det(\mathbf{E})=0$.
  \item 5 DOF (3 per $\mathbf{R}$, 2 per direzione di $\mathbf{t}$).
  \item 2 valori singolari non-zero uguali.
\end{itemize}
\end{propbox}

\begin{notebox}[E vs.\ H]
$\mathbf{E}$ mappa un punto in una linea ($\mathbf{l}'=\mathbf{Ex}$);
$\mathbf{H}$ mappa un punto in un punto ($\mathbf{x}'=\mathbf{Hx}$).
\end{notebox}

\section{La matrice fondamentale}

\begin{defbox}[Matrice fondamentale]
$\mathbf{F}\in\R^{3\times3}$ è la generalizzazione di $\mathbf{E}$ per camere
\textbf{non calibrate} ($\mathbf{K}$ ignota); opera direttamente su pixel:
\[
  \mathbf{p}'^T\mathbf{F}\mathbf{p}=0, \qquad
  \mathbf{F}=\mathbf{K}'^{-T}\mathbf{E}\mathbf{K}^{-1}
\]
\end{defbox}

\begin{propbox}[Derivazione algebrica di F da E]
Siano $\mathbf{p}, \mathbf{p}'$ le coordinate pixel e $\mathbf{x}, \mathbf{x}'$ le coordinate normalizzate della camera. La matrice intrinseca $\mathbf{K}$ lega i due domini:
\[
  \mathbf{p} = \mathbf{K} \mathbf{x} \implies \mathbf{x} = \mathbf{K}^{-1} \mathbf{p} \quad \text{e} \quad \mathbf{x}' = \mathbf{K}'^{-1} \mathbf{p}'
\]
Sostituendo nell'equazione del vincolo epipolare della matrice essenziale $\mathbf{x}'^T \mathbf{E} \mathbf{x} = 0$:
\[
  (\mathbf{K}'^{-1} \mathbf{p}')^T \mathbf{E} (\mathbf{K}^{-1} \mathbf{p}) = 0 \implies \mathbf{p}'^T \underbrace{(\mathbf{K}'^{-T} \mathbf{E} \mathbf{K}^{-1})}_{\mathbf{F}} \mathbf{p} = 0
\]
Otteniamo la definizione formale della Matrice Fondamentale:
\[
  \mathbf{F} = \mathbf{K}'^{-T} \mathbf{E} \mathbf{K}^{-1} = \mathbf{K}'^{-T} \mathbf{R} [\mathbf{t}]_\times \mathbf{K}^{-1}
\]
\end{propbox}

\begin{propbox}[Proprietà di $\mathbf{F}$]
\begin{itemize}
  \item Opera su coordinate pixel.
  \item Rango 2, $\det(\mathbf{F})=0$.
  \item 7 DOF.
  \item Fornisce una ``calibrazione debole''.
  \item $\mathbf{l}'=\mathbf{Fp}$.
\end{itemize}
\end{propbox}

\begin{notebox}[F vs.\ E vs.\ H]
$\mathbf{H}$: rango 3, 8 DOF. $\mathbf{F}$: rango 2, 7 DOF.
$\mathbf{E}$: rango 2, 5 DOF. Se $\mathbf{K},\mathbf{K}'$ note:
$\mathbf{E}=\mathbf{K}'^T\mathbf{F}\mathbf{K}$.
\end{notebox}

\section{Stima della matrice fondamentale}

\subsection{Algoritmo a 8 punti}

Ogni corrispondenza dà un'equazione lineare nelle 9 incognite di $\mathbf{F}$.
Con $\geq8$ corrispondenze: sistema $\mathbf{Af}=\mathbf{0}$
$\Rightarrow$ soluzione SVD (autovettore di $\mathbf{A}^T\mathbf{A}$
per il minimo autovalore).

\begin{warnbox}[Instabilità numerica]
Le colonne di $\mathbf{A}$ hanno magnitudini molto diverse
(es.\ $p_xp'_x$ vs.\ costante 1) $\Rightarrow$ cattivo condizionamento.
\end{warnbox}

\subsection{Algoritmo a 8 punti normalizzato}

\begin{enumerate}
  \item \textbf{Normalizzare}: traslare i punti al centroide; scalare a distanza
        media $\sqrt{2}$. Matrici: $\mathbf{T}$, $\mathbf{T}'$.
  \item \textbf{Stimare $\hat{\mathbf{F}}$} sui punti normalizzati (algoritmo
        a 8 punti base).
  \item \textbf{Imporre rango 2}: SVD di $\hat{\mathbf{F}}=\mathbf{U}\boldsymbol{\Sigma}\mathbf{V}^T$;
        imporre a zero il minimo valore singolare; ricalcolare $\hat{\mathbf{F}}'$.
  \item \textbf{Denormalizzare}: $\mathbf{F}=\mathbf{T}'^T\hat{\mathbf{F}}'\mathbf{T}$.
\end{enumerate}

In pratica si combina con \textbf{RANSAC} per gestire gli outlier.

\section{Triangolazione}

Dati $\mathbf{P},\mathbf{P}'$ e corrispondenze $(\mathbf{x},\mathbf{x}')$:

\[
  \mathbf{x}\times(\mathbf{PX})=\mathbf{0}, \qquad
  \mathbf{x}'\times(\mathbf{P'X})=\mathbf{0}
\]

Ogni prodotto vettoriale dà 2 equazioni in $\mathbf{X}\in\R^4$.
Sistema $4\times3$ $\mathbf{AX}=\mathbf{0}$ risolto per SVD.

\begin{warnbox}[Il problema del rumore]
I raggi retroproiettati non si intersecano perfettamente a causa del rumore.
Il metodo SVD minimizza un errore algebrico. Il \textbf{bundle adjustment}
minimizza l'errore di riproiezione (non lineare, più accurato).
\end{warnbox}

\section{Rettificazione stereo e DSI}

La \textbf{rettificazione stereo} deforma le due immagini per allineare le linee
epipolari orizzontalmente, riducendo il matching da 2D a 1D.

La \textbf{Disparity Space Image (DSI)} memorizza i costi di matching
$C(i,j)$ per ogni coppia di pixel corrispondenti nella stessa scanline.
La \textbf{programmazione dinamica} trova il percorso ottimale:
\[
  C(i,j)=\min\begin{cases}C(i-1,j-1)+D(i,j)\quad\text{(match)}\\C(i-1,j)+c_{\mathrm{occ}}\quad\text{(occlusione da destra)}\\C(i,j-1)+c_{\mathrm{occ}}\quad\text{(occlusione da sinistra)}\end{cases}
\]
Trattando le scanline indipendentemente si possono generare artefatti di striatura.

\section{Da $\mathbf{F}$ a moto della camera}

Se $\mathbf{K},\mathbf{K}'$ sono note, si recupera $\mathbf{E}=\mathbf{K}'^T\mathbf{FK}$
e si decompone: 4 soluzioni possibili per $(\mathbf{R},\mathbf{t})$.
La corretta si determina verificando che i punti 3D siano davanti a entrambe le
camere.

Per 3 viste: \textbf{tensore trifocale}; per 4 viste: \textbf{tensore quadrifocale}.

% =============================================================================
\chapter{Vision Transformers e Architetture Avanzate}
% =============================================================================

\section{Dal NLP alla Visione: Motivazione}

I \textbf{Transformers} sono stati introdotti originariamente per il Natural Language
Processing (NLP) nel paper \emph{"Attention Is All You Need"} (Vaswani et al., 2017).
La loro estensione alla visione artificiale deriva dalla capacità del meccanismo di
\textbf{Self-Attention} di modellare le interazioni tra coppie di elementi in una
sequenza in un'unica operazione globale, eliminando il bias di località tipico delle CNN.

% -- 12.2 Self-Attention -------------------------------------------------------
\section{Meccanismo di Self-Attention}

Dato un insieme di vettori di input $X \in \R^{n \times d_{model}}$, il modello impara
tre proiezioni lineari tramite matrici di peso $W_Q, W_K \in \R^{d_{model} \times d_k}$
e $W_V \in \R^{d_{model} \times d_v}$:
\[
  Q = X W_Q \quad (\text{Queries}), \qquad
  K = X W_K \quad (\text{Keys}), \qquad
  V = X W_V \quad (\text{Values})
\]

\begin{defbox}[Scaled Dot-Product Attention]
La risposta di attenzione viene calcolata moltiplicando la matrice delle Queries con
la trasposta delle Keys, applicando un fattore di scala e la funzione softmax sui
Values:
\[
  \mathrm{Attention}(Q, K, V) = \mathrm{softmax}\left(\frac{Q K^T}{\sqrt{d_k}}\right) V
\]
\end{defbox}

\begin{notebox}[Perché il fattore di scala $\frac{1}{\sqrt{d_k}}$?]
Per valori elevati della dimensione $d_k$, il prodotto scalare $Q K^T$ tende ad
assumere valori numericamente grandi. Questo spinge la funzione $\mathrm{softmax}$ in
regioni a gradiente estremamente piccolo (saturazione), provocando il fenomeno del
\emph{vanishing gradient} durante l'addestramento. La divisione per $\sqrt{d_k}$
stabilizza i gradienti mantenendo la varianza dei prodotti scalari vicina a $1$.
\end{notebox}

\subsection{Multi-Head Attention (MHA)}

Invece di calcolare un'unica funzione di attenzione, la \textbf{Multi-Head Attention}
esegue $h$ meccanismi di attenzione in parallelo. Ogni "testa" (head) proietta $Q, K, V$
in sottospazi di rappresentazione differenti, consentendo al modello di focalizzarsi
contemporaneamente su relazioni spaziali o semantiche diverse:
\[
  \mathrm{MultiHead}(Q, K, V) = \mathrm{Concat}(\text{head}_1, \ldots, \text{head}_h) W^O
\]
dove $\text{head}_i = \mathrm{Attention}(Q W_i^Q, K W_i^K, V W_i^V)$ e $W^O \in \R^{h d_v \times d_{model}}$.

\subsection{Positional Encoding}

La Self-Attention è \textbf{permutation-equivariant}: è del tutto indifferente
all'ordine d'ingresso dei token. Poiché nella visione artificiale la posizione
geometrica dei pixel è fondamentale, occorre iniettare l'informazione spaziale
sommando agli embedding di input un vettore di \textbf{Positional Encoding} (funzioni
sinusoidali fisse o embedding posizionali 1D/2D apprendibili).

\subsection{Struttura del Transformer Encoder Block}

Ogni blocco dell'encoder applica in sequenza:
\begin{enumerate}
  \item Multi-Head Self-Attention con connessione residuale e Layer Normalization.
  \item Feed-Forward Network (FFN) composta da due strati lineari con non-linearità
        (es.\ GELU o ReLU), integrata da connessione residuale e Layer Normalization.
\end{enumerate}

% -- 12.3 Vision Transformer (ViT) ---------------------------------------------
\section{Vision Transformer (ViT)}

Il \textbf{Vision Transformer} (Dosovitskiy et al., 2021) applica l'architettura del
Transformer Encoder direttamente a sequenze di patch d'immagine.

\begin{defbox}[Pipeline di ViT]
\begin{enumerate}
  \item \textbf{Patch Extraction}: un'immagine $H \times W \times C$ viene divisa in
        $N = \frac{HW}{P^2}$ patch non sovrapposti di dimensione $P \times P$ (es.\ $16 \times 16$).
  \item \textbf{Linear Embedding}: ogni patch viene appiattito in un vettore $P^2 C$ e
        proiettato linearmente in una dimensione $d_{model}$.
  \item \textbf{Class Token ([CLS])}: un token apprendibile $[CLS]$ viene aggiunto in
        testa alla sequenza; la sua rappresentazione finale all'uscita dell'encoder
        verrà usata per la classificazione globale.
  \item \textbf{Positional Encoding}: si sommano vettori posizionali 1D apprendibili.
  \item \textbf{Transformer Encoder}: la sequenza passa attraverso $L$ blocchi
        residuali MHA + FFN.
  \item \textbf{Classification Head}: un perceptron multistrato (MLP) viene applicato
        all'output del token $[CLS]$.
\end{enumerate}
\end{defbox}

\begin{warnbox}[CNN vs.\ ViT: Inductive Biases]
\begin{itemize}
  \item \textbf{CNN}: possiedono forti \emph{inductive biases} specifici per le immagini,
        ovvero la \textbf{Località} (pixel vicini sono correlati) e l'\textbf{Equivarianza alla Traslazione}
        (stessi filtri applicati ovunque). Imparano molto velocemente anche su piccoli dataset.
  \item \textbf{ViT}: hanno \emph{inductive biases} visivi molto \textbf{deboli}. Non assumono
        innatamente la geometria 2D dell'immagine.
  \item \textbf{Conseguenza}: i ViT richiedono pre-addestramento su dataset di scala
        massiva (es.\ ImageNet-21k o JFT-300M) per apprendere da soli le relazioni spaziali.
        Su grande scala superano nettamente le CNN modellando interazioni globali $O(N^2)$.
\end{itemize}
\end{warnbox}

% -- 12.4 Swin Transformer -----------------------------------------------------
\section{Swin Transformer: Shifted-Window Attention}

Il ViT classico calcola l'attenzione globale tra tutti i patch, con complessità
quadratica $O(N^2)$ rispetto alla risoluzione. Lo \textbf{Swin Transformer} (Liu et al.,
2021) introduce l'attenzione a finestre traslate per ottenere complessità lineare $O(N)$
e mappe di feature gerarchiche multi-scala.

\begin{propbox}[Meccanismi di Swin Transformer]
\begin{itemize}
  \item \textbf{Local Windows}: l'immagine viene divisa in finestre non sovrapposte
        $M \times M$ (es.\ $7 \times 7$ patch). La Self-Attention viene calcolata
        esclusivamente all'interno di ciascuna finestra $\Rightarrow$ Complessità $O(N)$.
  \item \textbf{Shifted Window Attention}: tra strati consecutivi, le finestre
        vengono traslate di $(\lfloor M/2 \rfloor, \lfloor M/2 \rfloor)$ pixel. Questo
        consente connessioni \emph{cross-window} tra patch adiacenti appartenenti a
        finestre diverse, mantenendo la complessità computazionale lineare.
  \item \textbf{Hierarchical Feature Maps}: accorpa patch adiacenti nei layer profondi
        costruendo una piramide di feature simile alle CNN, ideale per detection e segmentazione.
\end{itemize}
\end{propbox}

% -- 12.5 DETR -----------------------------------------------------------------
\section{DETR: End-to-End Object Detection with Transformers}

\textbf{DETR} (Carion et al., 2020) riformula l'Object Detection come un problema di
predizione diretta d'insieme (*set-prediction*), eliminando componenti euristiche
manuali come le \emph{anchor boxes} e la \emph{Non-Maximum Suppression (NMS)}.

\begin{defbox}[Architettura di DETR]
\begin{enumerate}
  \item \textbf{CNN Backbone}: estrae una mappa di caratteristiche dall'immagine
        (es.\ $H/32 \times W/32 \times C$).
  \item \textbf{Transformer Encoder}: elabora le feature spaziali appiattite insieme
        a positional encodings.
  \item \textbf{Transformer Decoder}: riceve $N_q$ \textbf{Object Queries} apprendibili
        (vettori posizionali trasformati che "interrogano" la scena) ed esegue
        Cross-Attention con l'output dell'encoder.
  \item \textbf{Prediction Heads}: MLP lineari mappano ciascuna query in una bounding
        box $(x, y, w, h)$ e una classe (incluso la classe "no object").
  \item \textbf{Bipartite Matching Loss}: usa l'\textbf{algoritmo Ungherese} durante
        l'addestramento per assegnare in modo unico ogni predizione a un oggetto
        ground-truth.
\end{enumerate}
\end{defbox}

% =============================================================================
\chapter{Modelli Generativi Profondi: GANs e Diffusione}
% =============================================================================

\section{Formulazione generale}

I modelli generativi profondi apprendono la distribuzione di probabilità sottostante
$p_{\text{data}}(x)$ a partire da dati ad alta dimensione (immagini) per campionare nuove
istanze sintetiche $x \sim p_{\theta}(x)$.

% -- 13.1 GANs -----------------------------------------------------------------
\section{Generative Adversarial Networks (GANs)}

Un'architettura \textbf{GAN} (Goodfellow et al., 2014) è composta da due reti neurali
addestrate in un gioco minimax avversario a due giocatori:
\begin{itemize}
  \item \textbf{Generatore ($G$)}: mappa un vettore latente stocastico $z \sim p_z$ in un
        campione sintetico $G(z)$.
  \item \textbf{Discriminatore ($D$)}: calcola la probabilità $D(x) \in [0, 1]$ che un
        campione provenga dai dati reali $p_{\text{data}}$ piuttosto che da $G$.
\end{itemize}

\begin{defbox}[Funzione Obiettivo Minimax delle GAN]
\[
  \min_G \max_D V(D, G) = \mathbb{E}_{x \sim p_{\text{data}}}[\log D(x)] + \mathbb{E}_{z \sim p_z}[\log(1 - D(G(z)))]
\]
All'equilibrio di Nash, $G$ ricostruisce esattamente la distribuzione $p_{\text{data}}$ e $D(x) \equiv 1/2$.
\end{defbox}

\begin{warnbox}[Problematiche di Addestramento e Fallimenti]
\begin{itemize}
  \item \textbf{Mode Collapse}: $G$ impara a produrre solo una piccola sotto-varietà
        di immagini reali che ingannano $D$, ignorando la diversità del dataset.
  \item \textbf{Vanishing Gradients}: se $D$ diventa troppo forte rispetto a $G$, i
        gradienti per $G$ crollano a zero.
  \item \textbf{Wasserstein GAN (WGAN)}: risolve le instabilità sostituendo la
        divergenza JS con la \emph{Distanza di Wasserstein-1} (Earth Mover's Distance) e
        imponendo il vincolo di $1$-Lipschitzianità a $D$.
\end{itemize}
\end{warnbox}

% -- 13.2 DDPM Diffusion Models ------------------------------------------------
\section{Denoising Diffusion Probabilistic Models (DDPM)}

I modelli di diffusione (Ho et al., 2020) definiscono un processo di generazione basato
su due catene stocastiche di Markov.

\subsection{Forward Process (Noising)}

Un processo fisso (non apprendibile) che corrompe progressivamente un'immagine reale $x_0$
aggiungendo rumore Gaussiano isotropo in $T$ passi temporali:
\[
  q(x_t \mid x_{t-1}) = \mathcal{N}\left(x_t; \sqrt{1 - \beta_t} x_{t-1}, \beta_t I\right)
\]
dove $\{\beta_t \in (0, 1)\}_{t=1}^T$ è una schedula di varianza prefissata.

\begin{propbox}[Reparameterization Trick nei DDPM]
Definendo $\alpha_t = 1 - \beta_t$ e $\bar{\alpha}_t = \prod_{i=1}^t \alpha_i$, è
possibile campionare direttamente l'immagine rumorosa $x_t$ al generico passo $t$ a partire
da $x_0$ in forma chiusa senza iterare tutti i passi precedenti:
\[
  x_t = \sqrt{\bar{\alpha}_t} x_0 + \sqrt{1 - \bar{\alpha}_t} \epsilon, \qquad \epsilon \sim \mathcal{N}(0, I)
\]
\end{propbox}

\subsection{Reverse Process (Denoising)}

Il processo inverso è una catena di Markov apprendibile parametrizzata da una rete
neurale $\epsilon_\theta(x_t, t)$ (solitamente un'architettura \textbf{U-Net}):
\[
  p_\theta(x_{t-1} \mid x_t) = \mathcal{N}\left(x_{t-1}; \mu_\theta(x_t, t), \Sigma_\theta(x_t, t)\right)
\]

\begin{defbox}[Loss Semplificata di DDPM]
Invece di prevedere l'immagine pulita, la rete viene addestrata a **prevedere il rumore
$\epsilon$** aggiunto al passo $t$:
\[
  \mathcal{L}_{\text{simple}} = \mathbb{E}_{x_0, t, \epsilon} \left[ \left\| \epsilon - \epsilon_\theta(x_t, t) \right\|^2 \right]
\]
\end{defbox}
\begin{itemize}
  \item \textbf{Diffusione}: addestramento stabile (nessun gioco avversario), evita il
        mode collapse e garantisce un'altissima diversità dei campioni.
  \item \textbf{Svantaggio}: campionamento iterativo lento ($T \approx 1000$ passaggi U-Net
        all'inferenza). Mitigato da metodi come \textbf{DDIM} ( deterministic sampling a $\sim 50$
        passi) e modelli di consistenza.
\end{itemize}
\end{notebox}

% =============================================================================
\chapter{Deep Learning per Video e Dati Temporali}
% =============================================================================

\section{Il Video come Segnale Spazio-Temporale}

Un video è un segnale 3D rappresentato da un tensore di dimensione $T \times H \times W \times C$,
dove $T$ è la dimensione temporale (sequenza di frame) e $H \times W$ sono le coordinate spaziali.

\section{Architetture per Video Recognition}

\begin{itemize}
  \item \textbf{Single-Frame Models}: applicano una CNN 2D indipendentemente su ogni
        frame e aggregano le predizioni tramite pooling (ignore l'ordine temporale).
  \item \textbf{2D CNN + RNN / LSTM}: estraggono feature da ciascun frame con una CNN
        2D e passano la sequenza a una rete ricorsiva per modellare le dipendenze temporali.
  \item \textbf{3D CNNs (es.\ C3D)}: estendono la convoluzione 2D aggiungendo una dimensione
        temporale ai kernel ($T_k \times K \times K$), elaborando congiuntamente spazio e tempo.
  \item \textbf{Two-Stream Networks}: dividono l'analisi in due rami separati:
        \begin{enumerate}
          \item *Spatial Stream*: elabora i frame RGB (apparenza).
          \item *Temporal Stream*: elabora stack di campi di \textbf{Optical Flow} (movimento).
        \end{enumerate}
\end{itemize}

% -- 14.1 TimeSformer ---------------------------------------------------------
\section{TimeSformer: Video Vision Transformer}

Estende il Vision Transformer al video trattando regioni spazio-temporali 3D come
\emph{tubelets}. Per evitare l'esplosione della memoria dovuta all'attenzione congiunta
spazio-tempo $O((T \times N)^2)$, introduce la \textbf{Divided Space-Time Attention}:

\begin{defbox}[Divided Space-Time Attention]
All'interno di ciascun blocco Transformer vengono applicati due meccanismi di attenzione
in sequenza:
\begin{enumerate}
  \item \textbf{Temporal Attention}: ogni patch interagisce esclusivamente con i patch
        situati nella stessa posizione spaziale attraverso tutti gli altri frame della sequenza.
  \item \textbf{Spatial Attention}: ogni patch interagisce con tutti gli altri patch
        all'interno dello stesso frame.
\end{enumerate}
Questo abbatte drasticamente i costi computazionali mantenendo la capacità di modellare
la dinamica del movimento.
\end{defbox}

% =============================================================================
\chapter{Vision-Language Models (VLMs)}
% =============================================================================

\section{Motivazione e Strategie di Fusione}

I modelli unimodali estraggono feature spaziali ma mancano di ragionamento semantico
linguistico. I \textbf{Vision-Language Models (VLMs)} rappresentano congiuntamente immagini
e testo per compiti quali \emph{Visual Question Answering (VQA)}, \emph{Image Captioning}
e \emph{Zero-Shot Recognition}.

\begin{propbox}[Strategie di Fusione Multimodale]
\begin{itemize}
  \item \textbf{Early Fusion}: concatenazione dei token visivi e testuali al primo strato di
        un unico Transformer condiviso.
  \item \textbf{Intermediate Fusion}: encoder dedicati separati le cui feature intermedie
        vengono fuse a metà rete tramite moduli di Cross-Attention (es.\ Flamingo).
  \item \textbf{Late Fusion}: elaborazione del tutto indipendente fino agli embedding
        finali, confrontati nello spazio latente (es.\ CLIP).
\end{itemize}
\end{propbox}

\begin{defbox}[Modality Gap]
Il \textbf{Modality Gap} indica la divergenza geometrica tra lo spazio delle feature visive
e quello delle feature testuali. Un VLM impara una proiezione condivisa in cui coppie
immagine-testo corrispondenti si trovino vicine nello spazio latente.
\end{defbox}

% -- 15.1 CLIP -----------------------------------------------------------------
\section{CLIP: Contrastive Language-Image Pre-training}

\textbf{CLIP} (Radford et al., 2021) è un VLM a Late Fusion addestrato su $400$ milioni di
coppie immagine-testo pubblicamente disponibili sul web tramite \textbf{Apprendimento Contrastivo}.

\begin{defbox}[InfoNCE Loss in CLIP]
Dato un batch di $N$ coppie (Immagine $I_k$, Testo $T_k$), si calcola la matrice di
similarità del coseno $N \times N$. La loss massimizza i valori lungo la diagonale
(coppie positive) e minimizza gli elementi fuori diagonale:
\[
  \mathcal{L}_{I \to T} = -\sum_{k=1}^N \log \frac{\exp(I_k \cdot T_k / \tau)}{\sum_{j=1}^N \exp(I_k \cdot T_j / \tau)}, \qquad
  \mathcal{L}_{\text{CLIP}} = \frac{1}{2}\left(\mathcal{L}_{I \to T} + \mathcal{L}_{T \to I}\right)
\]
dove $\tau$ è un parametro di temperatura apprendibile.
\end{defbox}

\subsection{Zero-Shot Transfer}

Durante l'inferenza, non occorre alcun addestramento specifico. I nomi delle classi
vengono trasformati in prompt testuali tramite *prompt engineering* (es.\ *"a photo of a \{class\}"*).
L'immagine di input viene classificata assegnandole la classe del prompt testuale il cui
embedding presenta la \textbf{massima similarità del coseno} con l'embedding dell'immagine.

% =============================================================================
\chapter{Medical Imaging e Segmentazione Avanzata}
% =============================================================================

\section{Il Medical Data Gap}

A differenza dei dataset visivi naturali (ImageNet), l'imaging medico presenta forti
limitazioni: volumetria ridotta, annotazioni estremamente costose poiché richiedono
medici specialisti, vincoli normativi rigidi sulla privacy (GDPR/HIPAA) e un forte
\textbf{Domain Shift} dovuto a scanner e protocolli di acquisizione differenti tra ospedali.

\section{Architetture di Segmentazione Biomedica}

\begin{defbox}[U-Net]
La \textbf{U-Net} (Ronneberger et al., 2015) è un'architettura Encoder-Decoder simmetrica:
\begin{itemize}
  \item \textbf{Encoder} (contrazione): estrae feature semantiche ad alto livello
        tramite convoluzioni e pooling.
  \item \textbf{Decoder} (espansione): ripristina la risoluzione spaziale originale
        tramite upsampling.
  \item \textbf{Skip Connections}: concatenano le mappe di caratteristiche ad alta
        risoluzione dell'encoder direttamente con i corrispondenti livelli del decoder,
        preservando i dettagli spaziali dei bordi delle lesioni.
\end{itemize}
\end{defbox}

\section{Funzioni di Loss per Dati Sbilanciati}

Nel settore medico, lo sfondo occupa la maggior parte dell'immagine, mentre il tumore/lesione
rappresenta solo una frazione minuscola di pixel.

\begin{propbox}[Dice Loss e Focal Loss]
\begin{itemize}
  \item \textbf{Dice Loss}: ottimizza direttamente l'overlap geometrico tra la maschera
        prevista $y$ e la ground-truth $t$, risultando insensibile allo sbilanciamento delle classi:
        \[
          \mathcal{L}_{\text{Dice}} = 1 - \frac{1}{C}\sum_{c=0}^{C-1} \frac{2 \sum_n t_n^c y_n^c}{\sum_n (t_n^c + y_n^c)}
        \]
  \item \textbf{Focal Loss}: abbassa il peso dei pixel facili (sfondo) focalizzando la
        backpropagation sugli esempi difficili (i confini della lesione):
        \[
          \mathcal{L}_{\text{Focal}} = -\sum_{n=1}^N (1 - p_{t,n})^\gamma \log(p_{t,n}), \qquad (\text{tipicamente } \gamma = 2)
        \]
\end{itemize}
\end{propbox}

\section{Metriche di Valutazione}

*   \textbf{Dice Score (DSC)}: equivalent all'F1-score $\implies \mathrm{DSC} = \frac{2|X \cap Y|}{|X| + |Y|}$.
*   \textbf{Hausdorff Distance (HD)}: misura la deviazione peggiore superficie-superficie tra
    i contorni previsti e quelli reali:
    \[
      \mathrm{HD}(X, Y) = \max \left\{ \max_{x \in X} \min_{y \in Y} \|x - y\|, \max_{y \in Y} \min_{x \in X} \|x - y\| \right\}
    \]

% =============================================================================
\chapter{Monocular Depth Estimation (MDE)}
% =============================================================================

\section{Definizione e Natura Ill-Posed}

La \textbf{Monocular Depth Estimation (MDE)} consiste nel prevedere una mappa di profondità
densa $Z(x,y)$ a partire da un'unica immagine RGB 2D.

\begin{warnbox}[Perché MDE è un problema mal posto?]
La proiezione da 3D a 2D perde un grado di libertà (la profondità lungo il raggio ottico).
Infiniti mondi 3D differenti possono proiettarsi nella stessa identica immagine 2D.
Il problema è quindi matematicamente mal posto e richiede che il modello apprenda
\emph{priors} visivi dal contesto.
\end{warnbox}

\section{Indizi Monoculari (Monocular Cues)}

*   \textbf{Prospettiva lineare}: linee parallele che convergono verso il punto di fuga.
*   \textbf{Dimensione relativa}: oggetti noti appaiono più piccoli se lontani.
*   \textbf{Gradiente di texture}: la grana delle superfici diventa più densa a distanza.
*   \textbf{Occlusione}: gli oggetti in primo piano nascondono quelli sullo sfondo.

\section{Depth Anything v2}

Modello state-of-the-art basato su una pipeline \textbf{Teacher-Student} di pseudo-labelling:
1. Si addestra un modello *Teacher* su dati sintetici (profondità perfetta senza rumore).
2. Il Teacher genera pseudo-labels dense su un dataset immenso di immagini reali non etichettate.
3. Si addestra un modello *Student* su queste pseudo-labels, garantendo forte generalizzazione.

% =============================================================================
\chapter{Efficienza nel Deep Learning}
% =============================================================================

\section{Training Efficiente vs.\ Inferenza Efficiente}

*   \textbf{Efficient Training}: riduce il tempo di addestramento su GPU/TPU (es.\ *mixed-precision FP16*, *gradient checkpointing*).
*   \textbf{Efficient Inference}: minimizza latenza e memoria durante la distribuzione su
    dispositivi mobile o embedded.

\section{Tecniche di Inferenza Efficiente}

\begin{defbox}[Tecniche di Compressione]
\begin{enumerate}
  \item \textbf{Architetture Compatte}: uso di \emph{depthwise separable convolutions} (es.\ MobileNet)
        per separare il filtraggio spaziale dalla combinazione dei canali.
  \item \textbf{Pruning}: rimozione dei pesi meno importanti.
        \begin{itemize}
          \item *Structured Pruning*: rimuove interi canali/filtri $\Rightarrow$ velocizzazione immediata su hardware standard.
          \item *Unstructured Pruning*: azzera pesi singoli $\Rightarrow$ richiede hardware specializzato per matrici sparse.
        \end{itemize}
  \item \textbf{Knowledge Distillation}: addestramento di una rete piccola (*Student*) affinché
        imiti la distribuzione delle probabilità umane o le feature intermedie di una rete grande (*Teacher*).
  \item \textbf{Quantizzazione}: riduzione della precisione numerica dei pesi e delle attivazioni.
        Passare da \textbf{FP32 ad INT8} riduce l'impronta di memoria di $4\times$ e abilita
        l'aritmetica intera sui processori.
\end{enumerate}
\end{defbox}

\section{Metriche e Trade-off}

*   \textbf{MACs}: operazioni di *Multiply-Accumulate*.
*   \textbf{FLOPs}: operazioni in virgola mobile ($\approx 2 \times \text{MACs}$).
*   \textbf{Pareto Frontier}: la curva che rappresenta le configurazioni di modelli ottimali
    per cui non esiste un'architettura contemporaneamente più precisa e più veloce.

\end{document}
